\documentclass[11pt]{article}

\usepackage{amsmath,amssymb,amsthm,amsfonts,mathtools,mathrsfs}
\usepackage{fullpage}
\usepackage[T1]{fontenc}
\usepackage{hyperref}
\usepackage{enumitem}
\usepackage{bm}

\setlength{\parskip}{0.45em}
\setlength{\parindent}{0pt}

\newtheorem{theorem}{Theorem}[section]
\newtheorem{lemma}[theorem]{Lemma}
\newtheorem{proposition}[theorem]{Proposition}
\newtheorem{corollary}[theorem]{Corollary}
\theoremstyle{definition}
\newtheorem{definition}[theorem]{Definition}
\newtheorem{hypothesis}[theorem]{Hypothesis}
\newtheorem{assumption}[theorem]{Assumption}
\theoremstyle{remark}
\newtheorem{remark}[theorem]{Remark}
\newtheorem{example}[theorem]{Example}

\DeclareMathOperator{\Res}{Res}
\DeclareMathOperator{\spec}{spec}
\DeclareMathOperator{\capa}{cap}
\DeclareMathOperator{\supp}{supp}

\title{The Riemann--Hilbert View of Pad\'e Approximation\\
and Global Analytic Continuation,\\
with Application to the Fractional Riccati Equation}
\author{}
\date{}

\begin{document}

\maketitle

\begin{abstract}
The diagonal Pad\'e approximant of a function analytic at the origin admits a representation as the solution to a $2\times 2$ matrix Riemann--Hilbert problem. This formulation replaces the local derivative-matching of Taylor approximation with the global reconstruction of a spectral measure from its moments. The convergence domain is determined not by the distance to the nearest singularity, but by the maximal region on which the associated Riemann--Hilbert factorization remains solvable, a region governed by the equilibrium measure of an extremal contour selected by the $S$-property. The poles of the approximant are the eigenvalues of a truncated Jacobi operator, and their condensation along the equilibrium contour is the signature of the underlying Riemann surface. We present the general theory in the form of theorems, corollaries, and limit statements, identify the precise hypotheses under which each conclusion holds, and apply the framework to the constant-coefficient fractional Riccati equation that governs the rough Heston characteristic function in mathematical finance. The application section develops the M\"untz spectral solution, derives a Gamma-ratio convolution recurrence for the Puiseux coefficients, establishes a strictly positive radius of convergence by a Banach fixed-point argument, constructs the diagonal Pad\'e resummation in the variable $z = t^{\mu}$, proves meromorphy of the transformed solution on $\mathbb{C}$ by a classical Riccati majorization with a sharp Gautschi-type constant and a $z_{*}$-dependent residue formula, and proves global convergence on the positive real axis, with a computable a-posteriori error bound.
\end{abstract}

\tableofcontents

\section{Introduction}
\label{sec:intro}

Taylor approximation asks for the polynomial that matches the derivatives of a function at a single point. It is local, and its radius of convergence is the distance to the nearest singularity. Pad\'e approximation asks for the spectral measure that generates the given moment sequence, and for the Stieltjes transform of that measure. It is global, and its convergence domain is the maximal region on which the Riemann--Hilbert factorization exists.

The distinction is not merely one of degree. A Taylor polynomial of degree $2n$ and a diagonal Pad\'e approximant $[n/n]$ both consume the same $2n+1$ coefficients, yet they encode fundamentally different information. The Taylor polynomial encodes derivatives at a point; the Pad\'e approximant encodes moments of a measure. The first is governed by local analyticity; the second by global potential theory. The convergence theory of the two paradigms reflects this distinction at every level.

The document proceeds as follows. Section~\ref{sec:rhp} presents the Fokas--Its--Kitaev Riemann--Hilbert formulation for diagonal Pad\'e approximants. Section~\ref{sec:stahl} describes Stahl's geometric theorem and the selection of the equilibrium contour by the $S$-property, together with the Deift--Zhou steepest descent asymptotic. Section~\ref{sec:froissart} explains why convergence holds in capacity rather than uniformly, through the mechanism of Froissart doublets, and gives the Gon\v{c}ar condition under which uniform convergence is recovered. Section~\ref{sec:jacobi} identifies the Pad\'e table with the Toda lattice and the Jacobi operator resolvent, including the non-Hermitian extension via pseudo-spectrum. Section~\ref{sec:poles} contrasts orthogonal polynomial zeros with Fekete points and gives bulk and edge residue estimates. Section~\ref{sec:borel} discusses the Borel plane and Pad\'e--Borel resummation for factorially divergent series. Section~\ref{sec:heston} develops the application to the rough Heston model in detail, treating the constant-coefficient fractional Riccati equation through the M\"untz spectral method and Pad\'e resummation in $z = t^{\mu}$. Section~\ref{sec:conclusion} summarizes the essential difference between the two paradigms.

\section{The Riemann--Hilbert formulation}
\label{sec:rhp}

We begin with the foundational equivalence between diagonal Pad\'e approximants and matrix Riemann--Hilbert problems. This equivalence, due to Fokas, Its, and Kitaev, recasts a problem of rational approximation as a problem of matrix factorization on a contour.

\begin{definition}[Cauchy transform]
\label{def:cauchy}
Let $\Sigma\subset\mathbb{C}$ be a finite union of smooth oriented arcs, and let $w:\Sigma\to\mathbb{C}$ be a measurable weight with $\int_{\Sigma}|w(s)|\,|ds|<\infty$. The \emph{Cauchy transform} of $w$ on $\Sigma$ is the function
\begin{equation}
f(z)=\int_{\Sigma}\frac{w(s)\,ds}{z-s},
\qquad z\in\mathbb{C}\setminus\Sigma.
\label{eq:cauchy}
\end{equation}
\end{definition}

\begin{definition}[Diagonal Pad\'e approximant]
\label{def:pade}
Let $f$ be analytic at infinity with formal expansion $f(z)=\sum_{k=0}^{\infty}a_{k}z^{-k-1}$. The \emph{diagonal Pad\'e approximant} $[n/n]_{f}$ of order $n$ is the unique rational function $P_{n}/Q_{n}$ with $\deg P_{n}<n$, $\deg Q_{n}=n$, and $Q_{n}$ monic, satisfying
\begin{equation}
f(z)-\frac{P_{n}(z)}{Q_{n}(z)}=O(z^{-2n-1})
\qquad\text{as }z\to\infty,
\label{eq:pade-defining}
\end{equation}
whenever the underlying Hankel determinant $\tau_{n}=\det(a_{i+j})_{0\le i,j\le n-1}$ is nonzero.
\end{definition}

\begin{theorem}[Fokas--Its--Kitaev representation]
\label{thm:fik}
Let $f$ be the Cauchy transform \eqref{eq:cauchy} of a weight $w$ on a fixed contour $\Sigma$, and assume the Hankel determinants $\tau_{1},\dots,\tau_{n}$ are nonzero. Then there exists a unique $2\times 2$ matrix-valued function $Y:\mathbb{C}\setminus\Sigma\to GL_{2}(\mathbb{C})$ satisfying:
\begin{enumerate}[label=(\roman*)]
\item $Y$ is analytic on $\mathbb{C}\setminus\Sigma$;
\item $Y$ has continuous boundary values $Y_{\pm}$ on $\Sigma$ related by the jump
\begin{equation}
Y_{+}(s)=Y_{-}(s)
\begin{pmatrix}
1 & 2\pi i\,w(s)\\
0 & 1
\end{pmatrix},
\qquad s\in\Sigma;
\label{eq:jump}
\end{equation}
\item $Y$ has the asymptotic behavior at infinity
\begin{equation}
Y(z)=\left(I+\frac{Y_{1}}{z}+\frac{Y_{2}}{z^{2}}+\cdots\right)
\begin{pmatrix}
z^{n} & 0\\
0 & z^{-n}
\end{pmatrix},
\qquad z\to\infty.
\label{eq:normalization}
\end{equation}
\end{enumerate}
The entries of $Y$ encode the diagonal Pad\'e approximant via
\begin{equation}
Y_{11}(z)=Q_{n}(z),
\qquad
Y_{12}(z)=\int_{\Sigma}\frac{Q_{n}(s)w(s)\,ds}{s-z},
\label{eq:y-entries}
\end{equation}
and the approximant itself is
\begin{equation}
[n/n]_{f}(z)=\frac{P_{n}(z)}{Q_{n}(z)}=\frac{Y_{12}(z)}{Y_{11}(z)}.
\label{eq:pade-from-y}
\end{equation}
\end{theorem}

\begin{corollary}[Orthogonality]
\label{cor:orthogonality}
The denominator $Q_{n}$ is the monic orthogonal polynomial of degree $n$ for the moment functional $\mathcal{L}[p]=\int_{\Sigma}p(s)w(s)\,ds$:
\begin{equation}
\int_{\Sigma}Q_{n}(s)s^{k}w(s)\,ds=0,
\qquad k=0,1,\dots,n-1.
\label{eq:orthogonality}
\end{equation}
\end{corollary}

\begin{remark}[Limitation: the contour must be given]
\label{rem:fixed-contour}
Theorem~\ref{thm:fik} requires a fixed contour $\Sigma$. For functions specified only by a power series, with branch cuts whose location is unknown, the contour is not given; it must be determined as part of the problem. The Riemann--Hilbert formulation in the form stated above becomes exact only after the equilibrium contour has been selected by the procedure of Section~\ref{sec:stahl}.
\end{remark}

\section{The contour is determined, not given}
\label{sec:stahl}

For functions whose singular structure is not specified in advance, the appropriate contour is selected by an extremal principle in logarithmic potential theory. This is the content of Stahl's geometric theorem.

\begin{definition}[Stahl class]
\label{def:stahl-class}
A function $f$ analytic at infinity belongs to the \emph{Stahl class} $\mathcal{S}$ if it admits multi-valued meromorphic continuation to a domain $\overline{\mathbb{C}}\setminus E$, where $E$ is a compact set of logarithmic capacity zero.
\end{definition}

\begin{definition}[Admissible compact]
\label{def:admissible}
A compact set $K\subset\mathbb{C}$ is \emph{admissible} for $f\in\mathcal{S}$ if $\overline{\mathbb{C}}\setminus K$ is a domain on which $f$ admits single-valued meromorphic continuation.
\end{definition}

\begin{theorem}[Stahl's geometric theorem]
\label{thm:stahl-geometric}
For every $f\in\mathcal{S}$, there exists a unique admissible compact $\Delta_{f}$ with the following properties:
\begin{enumerate}[label=(\roman*)]
\item $\overline{\mathbb{C}}\setminus\Delta_{f}$ is a connected domain $D_{f}$;
\item $f$ has single-valued meromorphic continuation to $D_{f}$;
\item $\capa(\Delta_{f})\le\capa(K)$ for every other admissible compact $K$.
\end{enumerate}
The set $\Delta_{f}$ is called the \emph{Stahl compact} or \emph{$S$-compact} of $f$.
\end{theorem}

\begin{theorem}[The $S$-property]
\label{thm:s-property}
Let $g(z,\infty)$ denote the Green's function of $D_{f}$ with pole at infinity. The arcs of $\Delta_{f}$ satisfy the symmetry condition
\begin{equation}
\frac{\partial g}{\partial n_{+}}(z)=\frac{\partial g}{\partial n_{-}}(z),
\qquad z\in\Delta_{f}\setminus\{\text{endpoints}\},
\label{eq:s-property}
\end{equation}
where $n_{\pm}$ are the two unit normals to the arc.
\end{theorem}

\begin{remark}[Variational interpretation]
\label{rem:variational}
The $S$-property is the Euler--Lagrange equation of a min-max variational problem. Among admissible compacts in a fixed homotopy class, the minimum of $\capa(K)$ is attained at $\Delta_{f}$, and the harmonic measure exerts equal pull on both sides of each arc.
\end{remark}

\begin{theorem}[Stahl's convergence theorem]
\label{thm:stahl-convergence}
Let $f\in\mathcal{S}$ with Stahl compact $\Delta_{f}$ and complementary domain $D_{f}$. The diagonal Pad\'e approximants converge to $f$ in capacity on compact subsets of $D_{f}\setminus\mathcal{P}_{f}$, where $\mathcal{P}_{f}$ is the pole set of $f$. For every compact $K\subset D_{f}\setminus\mathcal{P}_{f}$ and every $\varepsilon>0$,
\begin{equation}
\lim_{n\to\infty}\capa\bigl\{z\in K:\bigl|f(z)-[n/n]_{f}(z)\bigr|>\varepsilon\bigr\}=0.
\label{eq:convergence-in-capacity}
\end{equation}
\end{theorem}

\begin{theorem}[Deift--Zhou asymptotic]
\label{thm:deift-zhou}
Under the hypotheses of Theorem~\ref{thm:stahl-convergence}, with the additional regularity that $w$ is analytic in a neighborhood of $\Delta_{f}$ and the endpoints exhibit square-root behavior of the equilibrium density, nonlinear steepest descent yields
\begin{equation}
[n/n]_{f}(z)=f(z)\bigl(1+O(e^{-ng(z,\infty)})\bigr)
\label{eq:deift-zhou-asymptotic}
\end{equation}
uniformly on compact subsets of $D_{f}$ bounded away from $\Delta_{f}\cup\mathcal{P}_{f}$.
\end{theorem}

\begin{corollary}[Exponential rate of convergence]
\label{cor:exponential-rate}
\begin{equation}
\limsup_{n\to\infty}\bigl|f(z)-[n/n]_{f}(z)\bigr|^{1/n}\le e^{-g(z,\infty)},
\qquad z\in D_{f}\setminus\mathcal{P}_{f}.
\label{eq:rate}
\end{equation}
The level sets $\{g(z,\infty)=c\}$ are loci of equal convergence rate, and $\Delta_{f}=\{g=0\}$.
\end{corollary}

\begin{remark}[Comparison with Taylor]
\label{rem:taylor-comparison}
The Taylor polynomial of degree $2n$ converges geometrically only inside the disk of convergence, with rate $|z|/r$ where $r$ is the distance to the nearest singularity. The Pad\'e rate $e^{-g(z,\infty)}$ extends to all of $D_{f}$, generally a much larger domain. The two rates agree only inside the disk of convergence, where $g(z,\infty)\sim-\log|z|+\mathrm{const}$.
\end{remark}

\section{Froissart doublets and convergence in capacity}
\label{sec:froissart}

The convergence statement of Theorem~\ref{thm:stahl-convergence} is in capacity, not uniformly. The reason is the appearance of spurious poles, called Froissart doublets, that wander through the domain of holomorphy without condensing on the Stahl compact.

\begin{definition}[Froissart doublet]
\label{def:froissart}
A \emph{Froissart doublet} of $[n/n]_{f}$ is a pole $z_{j}^{(n)}\in D_{f}\setminus\mathcal{P}_{f}$ accompanied by a zero $\zeta_{j}^{(n)}$ at distance $|z_{j}^{(n)}-\zeta_{j}^{(n)}|=O(e^{-cn})$ for some $c>0$, with residue
\begin{equation}
\Res\bigl([n/n]_{f},z_{j}^{(n)}\bigr)=O(e^{-cn})
\qquad\text{as }n\to\infty.
\label{eq:doublet-residue}
\end{equation}
\end{definition}

\begin{proposition}[Boundedness of doublet count]
\label{prop:doublet-count}
For $f\in\mathcal{S}$, the number of Froissart doublets of $[n/n]_{f}$ in any compact $K\subset D_{f}\setminus\mathcal{P}_{f}$ is $O(1)$ as $n\to\infty$, but their location may oscillate erratically with $n$.
\end{proposition}

\begin{definition}[Gon\v{c}ar condition]
\label{def:goncar}
A function $f\in\mathcal{S}$ satisfies the \emph{Gon\v{c}ar condition} on a compact $K\subset D_{f}\setminus\mathcal{P}_{f}$ if the number of poles of $[n/n]_{f}$ in $K$ is uniformly bounded in $n$.
\end{definition}

\begin{theorem}[Gon\v{c}ar's uniform convergence theorem]
\label{thm:goncar}
If $f\in\mathcal{S}$ satisfies the Gon\v{c}ar condition on a compact $K\subset D_{f}\setminus\mathcal{P}_{f}$, then $[n/n]_{f}\to f$ uniformly on $K$. The exceptional set of spurious poles is empty for sufficiently large $n$.
\end{theorem}

\begin{remark}
\label{rem:goncar-restrictive}
The Gon\v{c}ar condition is restrictive for generic $f\in\mathcal{S}$ because Froissart doublets, though $O(1)$ in number, could in principle visit every compact infinitely often. For meromorphic functions with well-separated poles and a regular Pad\'e table, the condition holds on every compact disjoint from the pole set; this is the situation of Section~\ref{sec:heston-global}.
\end{remark}

\section{The Jacobi operator and the Toda lattice}
\label{sec:jacobi}

The Pad\'e problem admits a third equivalent formulation, in addition to rational approximation and Riemann--Hilbert factorization. This third formulation is operator-theoretic: the Pad\'e approximant is the resolvent of a truncated Jacobi operator.

\begin{definition}[Hankel matrix and tau function]
\label{def:hankel}
For a sequence $\{a_{k}\}_{k\ge 0}$, the \emph{Hankel matrix} of order $n$ is
\begin{equation}
H^{(n)}=\bigl(a_{i+j}\bigr)_{0\le i,j\le n-1},
\label{eq:hankel}
\end{equation}
and its determinant $\tau_{n}=\det H^{(n)}$ is the \emph{tau function}.
\end{definition}

\begin{theorem}[Toda lattice structure]
\label{thm:toda}
Assume $\tau_{n}\ne 0$ for all $n\ge 1$. The monic orthogonal polynomials $Q_{n}$ for the moment functional defined by $\{a_{k}\}$ satisfy the three-term recurrence
\begin{equation}
Q_{n+1}(z)=(z-b_{n})Q_{n}(z)-a_{n}^{2}Q_{n-1}(z),
\qquad Q_{-1}=0,\;Q_{0}=1.
\label{eq:three-term}
\end{equation}
Under the spectral evolution $d\mu_{t}(x)=e^{tx}\,d\mu(x)$, the recurrence coefficients evolve according to the Toda equations
\begin{equation}
\dot{b}_{n}=a_{n}^{2}-a_{n-1}^{2},
\qquad
\dot{a}_{n}=a_{n}(b_{n+1}-b_{n}).
\label{eq:toda}
\end{equation}
The flow is isospectral.
\end{theorem}

\begin{definition}[Jacobi operator]
\label{def:jacobi}
The \emph{infinite Jacobi operator} associated to $\{a_{n},b_{n}\}$ is the tridiagonal matrix
\begin{equation}
J_{\infty}=
\begin{pmatrix}
b_{0} & a_{0} & 0 & \cdots\\
a_{0} & b_{1} & a_{1} & \ddots\\
0 & a_{1} & b_{2} & \ddots\\
\vdots & \ddots & \ddots & \ddots
\end{pmatrix}
\label{eq:jacobi-matrix}
\end{equation}
acting on $\ell^{2}(\mathbb{N}_{0})$. Its $n\times n$ \emph{principal truncation} is $J_{n}=[J_{\infty}]_{0\le i,j\le n-1}$.
\end{definition}

\begin{theorem}[Resolvent representation, self-adjoint case]
\label{thm:resolvent-representation}
Assume the moment functional is positive, so that $J_{\infty}$ is self-adjoint and bounded with spectral measure $\mu$ supported on a compact subset of $\mathbb{R}$. Then
\begin{equation}
f(z)=\langle e_{0},(J_{\infty}-zI)^{-1}e_{0}\rangle
=\int\frac{d\mu(x)}{x-z},
\label{eq:resolvent}
\end{equation}
and
\begin{equation}
[n/n]_{f}(z)=\langle e_{0},(J_{n}-zI)^{-1}e_{0}\rangle
=\sum_{j=1}^{n}\frac{|\langle e_{0},\psi_{j}^{(n)}\rangle|^{2}}{\lambda_{j}^{(n)}-z},
\label{eq:truncated-resolvent}
\end{equation}
where $\{\lambda_{j}^{(n)},\psi_{j}^{(n)}\}_{j=1}^{n}$ are the eigenvalue-eigenvector pairs of $J_{n}$.
\end{theorem}

\begin{corollary}[Poles as eigenvalues]
\label{cor:poles-eigenvalues}
The poles of $[n/n]_{f}$ are precisely the eigenvalues of $J_{n}$, and the residues are the squared first components of the corresponding normalized eigenvectors.
\end{corollary}

\begin{theorem}[Strong resolvent convergence]
\label{thm:strong-resolvent}
In the setting of Theorem~\ref{thm:resolvent-representation}, $J_{n}\to J_{\infty}$ in the strong resolvent sense:
\begin{equation}
\lim_{n\to\infty}(J_{n}-zI)^{-1}\xi=(J_{\infty}-zI)^{-1}\xi
\qquad\forall\xi\in\ell^{2},\;z\notin\spec(J_{\infty}).
\label{eq:strong-resolvent}
\end{equation}
Consequently, $[n/n]_{f}(z)\to f(z)$ for every $z\in\mathbb{C}\setminus\spec(J_{\infty})$, locally uniformly.
\end{theorem}

\begin{corollary}[Self-adjoint convergence domain]
\label{cor:sa-domain}
For real measures, the convergence domain of the diagonal Pad\'e approximation coincides with $\mathbb{C}\setminus\spec(J_{\infty})$, and $\Delta_{f}=\spec(J_{\infty})$.
\end{corollary}

\begin{remark}[Non-Hermitian extension]
\label{rem:non-hermitian}
For complex measures and functions with branch points off the real axis, $J_{\infty}$ is non-Hermitian and Theorem~\ref{thm:strong-resolvent} does not apply directly. The relevant convergence theory uses the pseudo-spectrum
\begin{equation}
\spec_{\varepsilon}(J_{\infty})=\bigl\{z\in\mathbb{C}:\bigl\|(J_{\infty}-zI)^{-1}\bigr\|\ge\varepsilon^{-1}\bigr\}.
\label{eq:pseudospectrum}
\end{equation}
\end{remark}

\begin{theorem}[Pseudo-spectral limit, regular case]
\label{thm:pseudo-limit}
Suppose the recurrence coefficients $\{a_{n},b_{n}\}$ admit a regular asymptotic expansion in $1/n$ and the limiting density is supported on a finite union of analytic arcs. Then
\begin{equation}
\lim_{\varepsilon\to 0}\partial\spec_{\varepsilon}(J_{\infty})=\Delta_{f}.
\label{eq:pseudo-limit}
\end{equation}
\end{theorem}

\begin{remark}[Beyond the regular case]
\label{rem:non-regular}
In less regular non-Hermitian settings, the limiting object of \eqref{eq:pseudo-limit} is a quadratic differential trajectory in the sense of Strebel; its identification with $\Delta_{f}$ requires additional symmetry conditions and is established case-by-case.
\end{remark}

\section{Pole condensation: orthogonal polynomials and Fekete points}
\label{sec:poles}

\begin{definition}[Fekete points]
\label{def:fekete}
The \emph{$n$-th Fekete points} of a compact $K\subset\mathbb{C}$ are the maximizers
\begin{equation}
\{z_{1}^{F,n},\dots,z_{n}^{F,n}\}=\arg\max_{(z_{1},\dots,z_{n})\in K^{n}}\prod_{i<j}|z_{i}-z_{j}|.
\label{eq:fekete}
\end{equation}
\end{definition}

\begin{theorem}[Equilibrium limit of Fekete points]
\label{thm:fekete-equilibrium}
The normalized counting measures of the Fekete points of $K$ converge weak-$\ast$ to the equilibrium measure $\mu_{K}$:
\begin{equation}
\frac{1}{n}\sum_{j=1}^{n}\delta_{z_{j}^{F,n}}\xrightarrow{w^{\ast}}\mu_{K}
\qquad\text{as }n\to\infty.
\label{eq:fekete-limit}
\end{equation}
\end{theorem}

\begin{theorem}[Equilibrium limit of orthogonal polynomial zeros]
\label{thm:zeros-equilibrium}
For $f\in\mathcal{S}$ with Stahl compact $\Delta_{f}$, the normalized counting measures of the zeros of $Q_{n}$ converge weak-$\ast$ to the equilibrium measure $\mu_{\Delta_{f}}$:
\begin{equation}
\frac{1}{n}\sum_{j=1}^{n}\delta_{\zeta_{j}^{(n)}}\xrightarrow{w^{\ast}}\mu_{\Delta_{f}}.
\label{eq:zeros-limit}
\end{equation}
\end{theorem}

\begin{remark}[Distinct objects, common limit]
\label{rem:distinct-objects}
Theorems~\ref{thm:fekete-equilibrium} and \ref{thm:zeros-equilibrium} share a conclusion but describe different sequences. Fekete points are extremal for discrete logarithmic energy and have no recursive relationship across $n$. Orthogonal polynomial zeros are determined by the spectral truncation of the Jacobi operator and are linked across $n$ by the three-term recurrence \eqref{eq:three-term}. Only their normalized counting measures agree in the limit.
\end{remark}

\begin{theorem}[Bulk residue estimate]
\label{thm:residue-bulk}
Let $f\in\mathcal{S}$ satisfy the regularity conditions of Theorem~\ref{thm:deift-zhou}. For poles $z_{j}^{(n)}$ approaching the interior of $\Delta_{f}$,
\begin{equation}
\Res\bigl([n/n]_{f},z_{j}^{(n)}\bigr)=
\frac{P_{n}(z_{j}^{(n)})}{Q_{n}'(z_{j}^{(n)})}
=O\!\left(\frac{1}{n}\right),
\label{eq:residue-bulk}
\end{equation}
uniformly for poles bounded away from the endpoints.
\end{theorem}

\begin{theorem}[Edge residue scaling]
\label{thm:residue-edge}
For poles within an $O(n^{-2/3})$ neighborhood of an endpoint of $\Delta_{f}$, the rescaled residues satisfy
\begin{equation}
n^{1/3}\Res\bigl([n/n]_{f},z_{j}^{(n)}\bigr)\to\rho_{\mathrm{Ai}}(\xi_{j}),
\qquad\xi_{j}=n^{2/3}(z_{j}^{(n)}-\text{endpoint}),
\label{eq:residue-edge}
\end{equation}
where $\rho_{\mathrm{Ai}}$ is determined by the Airy kernel. The bulk $1/n$ law is replaced by an Airy-type $n^{-1/3}$ decay at the edge, the universal Painlev\'e II scaling of random matrix theory.
\end{theorem}

\begin{corollary}[Pole-zero cancellation on contours avoiding the cut]
\label{cor:cancellation}
For any rectifiable contour $\gamma\subset D_{f}\setminus\mathcal{P}_{f}$ enclosing a region disjoint from $\Delta_{f}$,
\begin{equation}
\lim_{n\to\infty}\frac{1}{2\pi i}\oint_{\gamma}[n/n]_{f}(z)\,dz
=\frac{1}{2\pi i}\oint_{\gamma}f(z)\,dz.
\label{eq:contour-integral}
\end{equation}
\end{corollary}

\section{The Borel plane and resummation}
\label{sec:borel}

\begin{definition}[Borel transform]
\label{def:borel}
For a formal power series $\hat{f}(z)=\sum_{n=0}^{\infty}a_{n}z^{n}$, the \emph{Borel transform} is
\begin{equation}
B(t)=\sum_{n=0}^{\infty}\frac{a_{n}}{n!}\,t^{n}.
\label{eq:borel-transform}
\end{equation}
\end{definition}

\begin{theorem}[Inverse Borel transform]
\label{thm:inverse-borel}
If $B(t)$ admits analytic continuation to a sector $|\arg t|<\theta$ with at most exponential growth, the Laplace integral
\begin{equation}
f_{\mathrm{Borel}}(z)=\int_{0}^{\infty}e^{-t/z}B(t)\,dt
\label{eq:inverse-borel}
\end{equation}
defines a holomorphic function in a sector of the $z$-plane.
\end{theorem}

\begin{theorem}[Pad\'e--Borel resummation]
\label{thm:pade-borel}
Let $\hat{f}$ have factorial coefficient growth $a_{n}\sim n!$, and assume $B(t)$ is in Stahl's class with no positive real singularities. The composition
\begin{equation}
f_{\mathrm{P-B}}(z)=
\int_{0}^{\infty}e^{-t/z}[n/n]_{B}(t)\,dt
\label{eq:pade-borel}
\end{equation}
converges in capacity to $f_{\mathrm{Borel}}(z)$ on the Borel domain of regularity, and assigns a finite value to a series whose original radius of convergence was zero.
\end{theorem}

\begin{remark}[Resummation versus continuation]
\label{rem:resummation}
Pad\'e--Borel resummation is distinct from analytic continuation: the latter extends a function across a singularity from a region where it already exists, while the former assigns a value to a divergent series. The Pad\'e approximant in \eqref{eq:pade-borel} approximates the measure in the Borel plane whose Laplace transform is the physical quantity.
\end{remark}

\section{Application to rough Heston: the constant-coefficient fractional Riccati equation}
\label{sec:heston}

The rough Heston characteristic function reduces, after standard transformations, to the constant-coefficient fractional Riccati equation
\begin{equation}
D^{\mu}y(t)=c_{0}+c_{1}y(t)+c_{2}y(t)^{2},
\qquad I^{1-\mu}y(0)=0,
\label{eq:fractional-riccati}
\end{equation}
where $D^{\mu}$ is the Caputo derivative of order $\mu\in(0,1]$, $H=\mu-\tfrac{1}{2}$ is the Hurst parameter, and the coefficients $c_{0},c_{1},c_{2}\in\mathbb{C}$ depend on the model parameters and the Fourier variable. We develop the M\"untz spectral solution, the Pad\'e resummation in the variable $z=t^{\mu}$, and the global convergence theorem with computable error bound.

\subsection{Fractional calculus and the M\"untz basis}
\label{sec:heston-prelim}

\begin{definition}[Riemann--Liouville integral and Caputo derivative]
\label{def:RL-Caputo}
For $r>0$ and $f\in L^{1}_{\mathrm{loc}}[0,T]$, the \emph{Riemann--Liouville integral} is
\begin{equation}
I^{r}f(t)=\frac{1}{\Gamma(r)}\int_{0}^{t}(t-s)^{r-1}f(s)\,ds.
\label{eq:RL-integral}
\end{equation}
For $\mu\in(0,1)$ and $f$ such that $I^{1-\mu}f\in C^{1}[0,T]$, the \emph{Caputo derivative} is
\begin{equation}
D^{\mu}f(t)=\frac{d}{dt}\bigl(I^{1-\mu}f(t)\bigr)-\frac{f(0)}{\Gamma(1-\mu)}t^{-\mu}.
\label{eq:caputo}
\end{equation}
The fundamental relations are $I^{\mu}D^{\mu}f(t)=f(t)-f(0)$ and $D^{\mu}I^{\mu}f(t)=f(t)$.
\end{definition}

\begin{proposition}[Action on power functions]
\label{prop:powers}
For $s>-1$ and $r>0$,
\begin{equation}
I^{r}t^{s}=\frac{\Gamma(s+1)}{\Gamma(s+r+1)}t^{s+r}.
\label{eq:Ir-power}
\end{equation}
For $\mu\in(0,1)$ and $s>0$,
\begin{equation}
D^{\mu}t^{s}=\frac{\Gamma(s+1)}{\Gamma(s+1-\mu)}t^{s-\mu}.
\label{eq:Dmu-power}
\end{equation}
\end{proposition}

\begin{proof}
For \eqref{eq:Ir-power}, the substitution $u=tx$ in the integral defining $I^{r}t^{s}$ yields a Beta function evaluation. For \eqref{eq:Dmu-power}, observe that $s>0$ implies $f(0)=0$, so the correction term in \eqref{eq:caputo} vanishes; differentiating $I^{1-\mu}t^{s}=\Gamma(s+1)/\Gamma(s+2-\mu)\cdot t^{s+1-\mu}$ and applying $\Gamma(s+2-\mu)=(s+1-\mu)\Gamma(s+1-\mu)$ gives the result.
\end{proof}

\begin{remark}[The Caputo derivative of a constant]
\label{rem:Dmu-const}
The hypothesis $s>0$ in \eqref{eq:Dmu-power} cannot be relaxed. For $s=0$, the correction term in \eqref{eq:caputo} cancels the formal power $t^{-\mu}/\Gamma(1-\mu)$ and yields $D^{\mu}(1)=0$, consistent with the convention that constants have zero Caputo derivative.
\end{remark}

\begin{definition}[M\"untz lattice]
\label{def:muntz}
Fix $\mu\in(0,1]$ and $T>0$. The \emph{M\"untz lattice} is
\begin{equation}
\mathcal{M}=\mathrm{span}\{1,t^{\mu},t^{2\mu},t^{3\mu},\ldots\}.
\label{eq:muntz-lattice}
\end{equation}
\end{definition}

\begin{theorem}[Closure of the M\"untz lattice]
\label{thm:closure}
For every $k\ge 1$ and $r\in\mu\mathbb{N}$:
\begin{align}
D^{\mu}t^{k\mu}&=\frac{\Gamma(k\mu+1)}{\Gamma((k-1)\mu+1)}t^{(k-1)\mu},
\label{eq:Dmu-muntz}\\
I^{r}t^{k\mu}&=\frac{\Gamma(k\mu+1)}{\Gamma(k\mu+r+1)}t^{k\mu+r},
\label{eq:Ir-muntz}\\
t^{j\mu}\cdot t^{k\mu}&=t^{(j+k)\mu}.
\label{eq:product-muntz}
\end{align}
Consequently $D^{\mu}\mathcal{M}\subset\mathcal{M}$, $I^{r}\mathcal{M}\subset\mathcal{M}$ for $r\in\mu\mathbb{N}$, and $\mathcal{M}$ is closed under multiplication.
\end{theorem}

\begin{remark}[Why the M\"untz lattice is the natural basis]
\label{rem:muntz-natural}
The three closure properties of Theorem~\ref{thm:closure}, under fractional differentiation, fractional integration of order $\mu$, and pointwise multiplication, make the M\"untz lattice the canonical basis for any constant-coefficient fractional ODE with polynomial nonlinearity. The M\"untz--Tau substitution preserves the algebraic structure of the equation. Ordinary Taylor series in $t$ fail because they are not closed under $D^{\mu}$ for non-integer $\mu$. Note that $I^{1-\mu}$ does \emph{not} preserve $\mathcal{M}$ for generic $\mu$; however, the Volterra reformulation \eqref{eq:volterra} requires only $I^{\mu}$, which does.
\end{remark}

\subsection{The series solution and Volterra representation}
\label{sec:heston-series}

Applying $I^{\mu}$ to \eqref{eq:fractional-riccati} and using $I^{\mu}D^{\mu}y=y-y(0)=y$ yields the equivalent Volterra integral equation
\begin{equation}
y(t)=\frac{1}{\Gamma(\mu)}\int_{0}^{t}(t-s)^{\mu-1}\bigl(c_{0}+c_{1}y(s)+c_{2}y(s)^{2}\bigr)\,ds.
\label{eq:volterra}
\end{equation}

\begin{remark}[H\"older regularity at the origin]
\label{rem:holder}
Evaluating \eqref{eq:volterra} at $t=0$ gives $y(0)=0$. Near the origin,
\begin{equation}
y(t)=\frac{c_{0}}{\Gamma(\mu+1)}t^{\mu}+O(t^{2\mu}),
\label{eq:holder-leading}
\end{equation}
so $y$ is H\"older continuous of order $\mu$ but not differentiable at $0$. The H\"older exponent $\mu$ is the natural regularity scale of the problem.
\end{remark}

\begin{theorem}[M\"untz--Tau Puiseux series solution]
\label{thm:riccati-series}
There exist $R>0$ and a sequence $(a_{k})_{k\ge 1}\subset\mathbb{C}$ such that
\begin{equation}
y(t)=\sum_{k=1}^{\infty}a_{k}t^{k\mu},
\qquad |t|<R^{1/\mu},
\label{eq:y-series}
\end{equation}
where the coefficients satisfy the Gamma-ratio convolution recurrence
\begin{align}
a_{1}&=\frac{c_{0}}{\Gamma(\mu+1)},
\label{eq:a1}\\
a_{k+1}&=\frac{\Gamma(k\mu+1)}{\Gamma((k+1)\mu+1)}\biggl(c_{1}a_{k}+c_{2}\!\!\sum_{\substack{j+\ell=k\\ 1\le j,\ell\le k-1}}\!\!a_{j}a_{\ell}\biggr),
\qquad k\ge 1.
\label{eq:ak}
\end{align}
\end{theorem}

\begin{proof}
Existence of $R>0$ is established in Theorem~\ref{thm:radius} below by a Banach fixed-point argument. Granting analyticity of $y$ in $z=t^{\mu}$ on $|z|<R$, write $y(t)=\sum_{k\ge 1}a_{k}t^{k\mu}$. By \eqref{eq:product-muntz},
\begin{equation}
y(t)^{2}=\sum_{m=2}^{\infty}b_{m}t^{m\mu},
\qquad b_{m}=\sum_{\substack{j+\ell=m\\ 1\le j,\ell\le m-1}}a_{j}a_{\ell}.
\label{eq:y-squared}
\end{equation}
Substituting into \eqref{eq:volterra}, splitting the integral, and applying \eqref{eq:Ir-power} with $s=k\mu$, $r=\mu$:
\begin{equation}
\int_{0}^{t}(t-s)^{\mu-1}s^{k\mu}\,ds=\Gamma(\mu)\frac{\Gamma(k\mu+1)}{\Gamma((k+1)\mu+1)}t^{(k+1)\mu}.
\label{eq:beta-integral}
\end{equation}
Matching the coefficient of $t^{n\mu}$ for each $n\ge 1$ yields \eqref{eq:a1} for $n=1$ and \eqref{eq:ak} after reindexing $n\mapsto k+1$.
\end{proof}

\begin{remark}[Quadratic convolution at low orders]
\label{rem:conv-low}
For $k=1$, the convolution sum in \eqref{eq:ak} is empty: there are no pairs $(j,\ell)$ with $j,\ell\ge 1$ and $j+\ell=1$. Hence
\begin{equation}
a_{2}=\frac{\Gamma(\mu+1)}{\Gamma(2\mu+1)}c_{1}a_{1}.
\label{eq:a2}
\end{equation}
The quadratic term first contributes at $k=2$:
\begin{equation}
a_{3}=\frac{\Gamma(2\mu+1)}{\Gamma(3\mu+1)}\bigl(c_{1}a_{2}+c_{2}a_{1}^{2}\bigr).
\label{eq:a3}
\end{equation}
\end{remark}

\subsection{Radius of convergence and generic finiteness}
\label{sec:heston-radius}

\begin{theorem}[Strictly positive Puiseux radius and Cauchy bound]
\label{thm:radius}
Let
\begin{equation}
A=|c_{0}|+|c_{1}|+|c_{2}|,\qquad B=|c_{1}|+2|c_{2}|,
\label{eq:AB}
\end{equation}
and define
\begin{equation}
T_{0}^{\mu}=\min\!\left(\frac{\Gamma(\mu+1)}{2A},\;\frac{\Gamma(\mu+1)}{4B}\right)>0.
\label{eq:T0}
\end{equation}
\begin{enumerate}[label=(\roman*)]
\item There exists a unique solution $y\in C[0,T_{0}]$ of \eqref{eq:volterra} with $\|y\|_{\infty}\le 1$. This solution is analytic in $z=t^{\mu}$ on $\{|z|<T_{0}^{\mu}\}$, so the Puiseux series \eqref{eq:y-series} has radius of convergence $R\ge T_{0}^{\mu}>0$.
\item For every $0<r<R$, with $g(z)=y(z^{1/\mu})$ and $M(r)=\sup_{|z|=r}|g(z)|$,
\begin{equation}
|a_{k}|\le \frac{M(r)}{r^{k}}\qquad \forall k\ge 1.
\label{eq:cauchy-bound}
\end{equation}
\end{enumerate}
\end{theorem}

\begin{proof}
Define $B_{1}=\{y\in C[0,T_{0}]:\|y\|_{\infty}\le 1\}$ with the supremum norm and the operator
\begin{equation}
\Phi(y)(t)=\frac{1}{\Gamma(\mu)}\int_{0}^{t}(t-s)^{\mu-1}\bigl(c_{0}+c_{1}y(s)+c_{2}y(s)^{2}\bigr)\,ds=I^{\mu}(c_{0}+c_{1}y+c_{2}y^{2})(t).
\label{eq:phi-op}
\end{equation}
For $y\in B_{1}$, the integrand is bounded by $A$, and \eqref{eq:Ir-power} gives
\begin{equation}
|\Phi(y)(t)|\le A\cdot\frac{T_{0}^{\mu}}{\Gamma(\mu+1)}\le\frac{1}{2}<1
\label{eq:phi-bound}
\end{equation}
by the first term in \eqref{eq:T0}. For $y_{1},y_{2}\in B_{1}$,
\begin{equation}
|c_{1}(y_{1}-y_{2})+c_{2}(y_{1}^{2}-y_{2}^{2})|\le B\|y_{1}-y_{2}\|_{\infty},
\label{eq:lipschitz}
\end{equation}
so
\begin{equation}
\|\Phi(y_{1})-\Phi(y_{2})\|_{\infty}\le \frac{T_{0}^{\mu}B}{\Gamma(\mu+1)}\|y_{1}-y_{2}\|_{\infty}\le\frac{1}{4}\|y_{1}-y_{2}\|_{\infty}<\frac{1}{2}\|y_{1}-y_{2}\|_{\infty}
\label{eq:contraction}
\end{equation}
by the second term in \eqref{eq:T0}. The Banach fixed-point theorem provides a unique $y^{\ast}\in B_{1}$ solving \eqref{eq:volterra}. The formal Puiseux series defined by \eqref{eq:a1}--\eqref{eq:ak} solves the same equation order by order; by uniqueness, the series equals $y^{\ast}$ on $[0,T_{0}]$, proving (i). Statement (ii) is Cauchy's integral formula applied to the analytic function $g$.
\end{proof}

\begin{theorem}[Generic finiteness of the radius]
\label{thm:radius-finite}
For generic parameters with $c_{2}\ne 0$, the Puiseux radius is finite: $R<\infty$. The complexified Riccati flow develops movable singularities at finite distance from the origin in the complex $z$-plane.
\end{theorem}

\begin{proof}[Sketch]
For $\mu=1$ with $c_{2}\ne 0$, the substitution $y=-w'/(c_{2}w)$ reduces $y'=c_{0}+c_{1}y+c_{2}y^{2}$ to a second-order linear ODE with constant coefficients for $w$. The zeros of $w$ generically lie at finite $t$, and at each zero $y$ has a simple pole. The algebraic structure of the recurrence \eqref{eq:ak} persists for $\mu\in(0,1)$, and the resulting $g(z)$ generically inherits movable singularities at finite distance. The exact location of the nearest singularity is encoded by the limiting distribution of the diagonal Pad\'e poles (Theorem~\ref{thm:zeros-equilibrium}).
\end{proof}

\begin{remark}[The recurrence is only a local representation]
\label{rem:local}
Theorem~\ref{thm:radius-finite} is the central obstruction. Even though every $a_{k}$ is computed in closed form by \eqref{eq:ak}, the series \eqref{eq:y-series} does not converge for all $t\ge 0$. The complex $z$-plane carries movable singularities of the Riccati flow, and $R$ is the distance to the nearest such singularity. On the positive real $t$-axis the solution $y(t)$ may exist for all $t$, yet the M\"untz series ceases to converge once $t^{\mu}$ exceeds $R$. A genuine global representation must exploit \emph{analytic continuation}, not raw series summation. This is the role played by Pad\'e approximation in the next subsection.
\end{remark}

\subsection{Pad\'e resummation in the variable $z=t^{\mu}$}
\label{sec:heston-pade}

Let $z=t^{\mu}$ and define
\begin{equation}
g(z)=y(z^{1/\mu})=\sum_{k=1}^{\infty}a_{k}z^{k},
\qquad |z|<R.
\label{eq:g-series}
\end{equation}
Fix $M\in\mathbb{N}$ and seek polynomials
\begin{equation}
P_{M}(z)=\sum_{k=0}^{M}p_{k}z^{k},
\qquad Q_{M}(z)=1+\sum_{k=1}^{M}q_{k}z^{k}
\label{eq:PQ-def}
\end{equation}
satisfying the matching condition
\begin{equation}
Q_{M}(z)g(z)-P_{M}(z)=O(z^{2M+1})
\qquad (z\to 0).
\label{eq:pade-match}
\end{equation}

\begin{lemma}[Hankel form of the denominator system]
\label{lem:hankel-form}
The matching condition \eqref{eq:pade-match} is equivalent to $p_{0}=0$,
\begin{align}
p_{n}&=a_{n}+\sum_{j=1}^{\min(n,M)}q_{j}a_{n-j},
\qquad 1\le n\le M,
\label{eq:p-eqs}\\
0&=a_{n}+\sum_{j=1}^{M}q_{j}a_{n-j},
\qquad M+1\le n\le 2M.
\label{eq:q-eqs}
\end{align}
The denominator system \eqref{eq:q-eqs} is the Hankel linear system
\begin{equation}
H_{M}\bm{q}=-\bm{b},
\qquad
H_{M}=
\begin{pmatrix}
a_{M} & a_{M-1} & \cdots & a_{1}\\
a_{M+1} & a_{M} & \cdots & a_{2}\\
\vdots & \vdots & \ddots & \vdots\\
a_{2M-1} & a_{2M-2} & \cdots & a_{M}
\end{pmatrix},
\quad
\bm{q}=
\begin{pmatrix}q_{1}\\ \vdots\\ q_{M}\end{pmatrix},
\quad
\bm{b}=
\begin{pmatrix}a_{M+1}\\ \vdots\\ a_{2M}\end{pmatrix}.
\label{eq:hankel-system}
\end{equation}
\end{lemma}

\begin{proof}
Direct expansion of $Q_{M}\cdot g$ and identification of coefficients of $z^{n}$ for $0\le n\le 2M$.
\end{proof}

\begin{proposition}[Regularity of the Pad\'e table]
\label{prop:regularity}
For $\mu\in(0,1)$, the function $g(z)=y(z^{1/\mu})$ is not rational: the fractional Riccati flow generates branch points that prevent any finite-order rational representation. Consequently $\det H_{M}\ne 0$ for every $M\ge 1$, and the diagonal Pad\'e table is regular.

For $\mu=1$, the Cole--Hopf linearization $y=-w'/(c_{2}w)$ reduces the equation to a second-order linear ODE with constant coefficients; the solution is meromorphic in $\mathbb{C}$. The Hankel matrix $H_{M}$ is singular only in the degenerate case that the solution happens to be rational of degree $<M$. These rational configurations are the exceptional points; for generic parameters, $\det H_{M}\ne 0$ for all $M\ge 1$. If a particular order yields $\det H_{M}=0$, the index $M\pm 1$ may be used.
\end{proposition}

\begin{lemma}[Explicit solution]
\label{lem:explicit}
If $\det H_{M}\ne 0$, the system \eqref{eq:hankel-system} has the unique solution $\bm{q}=-H_{M}^{-1}\bm{b}$, and the numerator coefficients are recovered from \eqref{eq:p-eqs}.
\end{lemma}

\begin{definition}[Diagonal Pad\'e approximant and its domain]
\label{def:pade-domain}
The diagonal Pad\'e approximant of order $M$ is
\begin{equation}
R_{M}(z)=\frac{P_{M}(z)}{Q_{M}(z)}=\frac{\sum_{k=0}^{M}p_{k}z^{k}}{1+\sum_{k=1}^{M}q_{k}z^{k}},
\label{eq:RM}
\end{equation}
and its \emph{Pad\'e domain} is
\begin{equation}
\mathcal{D}_{M}=\{z\in\mathbb{C}:\det H_{M}\ne 0,\;Q_{M}(z)\ne 0\}.
\label{eq:DM}
\end{equation}
\end{definition}

\subsection{Remainder representation and the computable error bound}
\label{sec:heston-error}

\begin{theorem}[Structural remainder representation]
\label{thm:remainder-structure}
Define $\varepsilon_{M}(z)=g(z)-R_{M}(z)$. Then
\begin{equation}
\varepsilon_{M}(z)=\frac{\Pi_{M+1}(z)}{Q_{M}(z)^{2}}z^{2M+1}
\label{eq:remainder-structure}
\end{equation}
for some polynomial $\Pi_{M+1}$ of degree at most $M+1$. In particular $\varepsilon_{M}(z)=O(z^{2M+1})$ as $z\to 0$.
\end{theorem}

\begin{definition}[Successive Pad\'e differences]
\label{def:differences}
For $k\ge 1$, the \emph{successive Pad\'e difference} is
\begin{equation}
\Delta_{k}(z)=R_{k}(z)-R_{k-1}(z),
\qquad R_{0}\equiv 0.
\label{eq:Delta}
\end{equation}
\end{definition}

\begin{assumption}[Uniform geometric tail]
\label{ass:geometric-tail}
At a fixed $z\in\mathcal{D}_{M}\cap\mathcal{D}_{M-1}$:
\begin{enumerate}[label=(\roman*)]
\item $g$ has no poles in $\{|\zeta|\le|z|\}$;
\item $Q_{M}(z)\ne 0$ and $Q_{M-1}(z)\ne 0$;
\item there exists $\rho\in(0,1)$ such that $|\Delta_{k+1}(z)|\le\rho|\Delta_{k}(z)|$ for every $k\ge M-1$.
\end{enumerate}
\end{assumption}

\begin{theorem}[Computable a-posteriori error bound]
\label{thm:error-bound}
Under Assumption~\ref{ass:geometric-tail} with contraction ratio $\rho$, the Pad\'e sequence $R_{k}(z)$ converges, and
\begin{equation}
|g(z)-R_{M}(z)|\le \frac{\rho|\Delta_{M}(z)|}{1-\rho}.
\label{eq:error-uniform}
\end{equation}
Defining the observable contraction ratio
\begin{equation}
\rho_{\ast}=\frac{|\Delta_{M}(z)|}{|\Delta_{M-1}(z)|}\le\rho,
\label{eq:rho-star}
\end{equation}
and assuming $|\Delta_{M}(z)|<|\Delta_{M-1}(z)|$ so that $\rho_{\ast}\in(0,1)$, we obtain
\begin{equation}
\boxed{\;|g(z)-R_{M}(z)|\le \frac{|\Delta_{M}(z)|^{2}}{|\Delta_{M-1}(z)|-|\Delta_{M}(z)|}.\;}
\label{eq:error-final}
\end{equation}
\end{theorem}

\begin{proof}
By Assumption~\ref{ass:geometric-tail}(iii) and induction, $|\Delta_{M+j}(z)|\le\rho^{j}|\Delta_{M}(z)|$ for every $j\ge 0$. The telescoping series
\begin{equation}
g(z)-R_{M}(z)=\sum_{k=M+1}^{\infty}\Delta_{k}(z)
\label{eq:telescoping}
\end{equation}
converges absolutely and is bounded by $|\Delta_{M}(z)|\sum_{j=1}^{\infty}\rho^{j}=\rho|\Delta_{M}(z)|/(1-\rho)$, which is \eqref{eq:error-uniform}. Substituting $\rho=\rho_{\ast}$ and simplifying yields \eqref{eq:error-final}.
\end{proof}

\begin{corollary}[Practical convergence certificate]
\label{cor:practical-cert}
The bound \eqref{eq:error-final} depends only on the two most recent Pad\'e iterates and is computable in $O(1)$ time given $R_{M-1}$ and $R_{M}$. It serves as an a-posteriori certificate: when $|\Delta_{M}|^{2}/(|\Delta_{M-1}|-|\Delta_{M}|)<2^{-\mathrm{bits}}$, the Pad\'e iterate $R_{M}(z)$ approximates $g(z)$ to the requested precision.
\end{corollary}

\subsection{Meromorphy on $\mathbb{C}$ via classical Riccati majorization}
\label{sec:heston-meromorphy}

\begin{theorem}[Meromorphy of the transformed solution]
\label{prop:meromorphy}
Let $c_{0},c_{1},c_{2}\in\mathbb{C}$ with $c_{2}\ne 0$ and $\mu\in(0,1]$. The function $g(z)=y(z^{1/\mu})$ defined by the M\"untz series \eqref{eq:y-series} extends to a meromorphic function on $\mathbb{C}$ with discrete pole set $\mathcal{P}_{g}$ having no finite accumulation point. Each pole $z_{*}\in\mathcal{P}_{g}$ is simple with residue
\begin{equation}
\Res(g,z_{*})=-\frac{(z_{*})^{(\mu-1)/\mu}}{c_{2}\,\mu}.
\label{eq:residue}
\end{equation}
\end{theorem}

\begin{proof}
The proof has six steps.

\emph{Step 1: sharp bound on the Gamma ratio.} Define $\gamma_{k}:=\Gamma(k\mu+1)/\Gamma((k+1)\mu+1)$ for $k\ge 1$, $\mu\in(0,1]$.

\begin{quote}
\textsc{Lemma A.} $C_{\mu}:=\sup_{k\ge 1}\gamma_{k}=\gamma_{1}=\Gamma(\mu+1)/\Gamma(2\mu+1)\le 21/20$ for all $\mu\in(0,1]$.
\end{quote}

\emph{Proof of Lemma A.} The log-convexity of $\Gamma$ on $(0,\infty)$ implies that for any fixed $\mu>0$ the function $k\mapsto\Gamma(k\mu+1)/\Gamma((k+1)\mu+1)=B(k\mu+1,\mu)/\Gamma(\mu)$ is strictly decreasing in $k$, because by the Beta-function identity $B(x+1,\mu)/B(x,\mu)=x/(x+\mu)<1$. Hence $\gamma_{k+1}/\gamma_{k}=(k\mu+1)/((k+1)\mu+1)<1$, so the supremum is attained at $k=1$. For the numerical bound at $k=1$, define $h(\mu)=\Gamma(\mu+1)/\Gamma(2\mu+1)$ on $[0,1]$. Then $h(0)=1$, $h(1)=1/2$, and $h'(\mu)=h(\mu)\bigl(\psi(\mu+1)-2\psi(2\mu+1)\bigr)$. Since $\psi(1)=-\gamma_{\mathrm{EM}}\approx-0.577$, we have $h'(0)=\psi(1)-2\psi(1)=-\psi(1)=\gamma_{\mathrm{EM}}>0$, so $h$ initially increases from $1$. The maximum of $h$ on $[0,1]$ is attained at the interior critical point where $\psi(\mu+1)=2\psi(2\mu+1)$; numerical solution gives $\mu_{*}\approx 0.1339$ with $h(\mu_{*})\approx 1.0428$. Hence
\begin{equation}
C_{\mu}\le\max_{\mu\in(0,1]}h(\mu)\le 1.0428<\frac{21}{20}.
\label{eq:Cmu-bound}
\end{equation}
(An earlier version of this argument claimed $C_{\mu}\le 1$; this was wrong precisely because $h$ exceeds $1$ near $\mu_{*}\approx 0.1339$. Every earlier estimate that used ``$\gamma_{k}\le 1$'' is hereby corrected.) In particular $\gamma_{k}\le C_{\mu}\le 21/20$ uniformly in $k\ge 1$ and $\mu\in(0,1]$.

\emph{Step 2: classical Riccati meromorphic majorant.} Define the majorant coefficients $\tilde A_{k}$ by
\begin{equation}
\tilde A_{1}=\frac{|c_{0}|}{\Gamma(\mu+1)},
\qquad
\tilde A_{k+1}=C_{\mu}\biggl(|c_{1}|\tilde A_{k}+|c_{2}|\!\!\sum_{j+\ell=k}\!\!\tilde A_{j}\tilde A_{\ell}\biggr),
\quad k\ge 1,
\label{eq:majorant-coeff}
\end{equation}
with the constant $C_{\mu}$ in front of \emph{every} nonlinear term. These are the power-series coefficients of the solution $\tilde y$ of the classical Riccati initial-value problem
\begin{equation}
\tilde y'(x)=\alpha+\beta\,\tilde y(x)+\gamma\,\tilde y(x)^{2},
\qquad \tilde y(0)=0,
\label{eq:majorant-ode}
\end{equation}
with $\alpha=C_{\mu}|c_{0}|/\Gamma(\mu+1)$, $\beta=C_{\mu}|c_{1}|$, $\gamma=C_{\mu}|c_{2}|$. The explicit solution is the classical Riccati quotient
\begin{equation}
\tilde y(x)=\frac{2\alpha}{\hat\beta}\tan\Bigl(\frac{\hat\beta}{2}x+\varphi_{0}\Bigr)-\frac{\beta}{2\gamma},
\qquad \hat\beta=\sqrt{4\alpha\gamma-\beta^{2}},\quad \varphi_{0}=\arctan\frac{\beta}{\hat\beta},
\label{eq:majorant-tan}
\end{equation}
if $\beta^{2}-4\alpha\gamma<0$, and
\begin{equation}
\tilde y(x)=\frac{\lambda_{+}\lambda_{-}\bigl(e^{(\lambda_{+}-\lambda_{-})x}-1\bigr)}{\lambda_{+}e^{(\lambda_{+}-\lambda_{-})x}-\lambda_{-}},
\qquad \lambda_{\pm}=\frac{\beta\pm\sqrt{\beta^{2}+4\alpha\gamma}}{2},
\label{eq:majorant-exp}
\end{equation}
if $\beta^{2}+4\alpha\gamma>0$ and $\alpha\gamma>0$ (the generic case), with the $\lambda_{\pm}$ scaled by $C_{\mu}$ through $\alpha,\beta,\gamma$. In either case $\tilde y$ is a \emph{meromorphic} function on $\mathbb{C}$ with a \emph{discrete} set $\widetilde{\mathcal{P}}$ of \emph{simple poles} with no finite accumulation point.

\emph{Step 3: coefficient majorization.} We claim $|a_{k}|\le\tilde A_{k}$ for all $k\ge 1$. The base case is $|a_{1}|=|c_{0}|/\Gamma(\mu+1)=\tilde A_{1}$ by \eqref{eq:a1}. For the inductive step, assuming $|a_{j}|\le\tilde A_{j}$ for all $j\le k$, the recurrence \eqref{eq:ak} gives
\begin{equation}
|a_{k+1}|
=\Bigl|\gamma_{k}\Bigl(c_{1}a_{k}+c_{2}\sum_{j+\ell=k}a_{j}a_{\ell}\Bigr)\Bigr|
\le\gamma_{k}\Bigl(|c_{1}||a_{k}|+|c_{2}|\sum_{j+\ell=k}|a_{j}||a_{\ell}|\Bigr)
\le C_{\mu}\Bigl(|c_{1}|\tilde A_{k}+|c_{2}|\sum_{j+\ell=k}\tilde A_{j}\tilde A_{\ell}\Bigr)
=\tilde A_{k+1},
\label{eq:majorization-step}
\end{equation}
using $\gamma_{k}\le C_{\mu}$ from Step 1. Therefore $\limsup_{k\to\infty}|a_{k}|^{1/k}\le 1/\tilde R$, where $\tilde R>0$ is the distance from $0$ to the nearest pole of $\tilde y$. Hence $g(z)$ converges on $|z|<\tilde R$ and on this disk
\begin{equation}
|g(z)|\le\sum_{k\ge 1}|a_{k}||z|^{k}\le\sum_{k\ge 1}\tilde A_{k}|z|^{k}=\tilde y(|z|).
\label{eq:majorant-bound}
\end{equation}

\emph{Step 4: tube continuation through $\mathbb{C}\setminus\widetilde{\mathcal{P}}$.} Fix $z_{0}\in\mathbb{C}\setminus\widetilde{\mathcal{P}}$ and let $\gamma\colon[0,1]\to\mathbb{C}\setminus\widetilde{\mathcal{P}}$ be a $C^{1}$ path with $\gamma(0)=0$, $\gamma(1)=z_{0}$. Define $d_{0}:=\operatorname{dist}\bigl(\gamma([0,1]),\widetilde{\mathcal{P}}\bigr)>0$. For $\varepsilon\in(0,d_{0}/2)$, let $\mathrm{Tube}_{\varepsilon}:=\{z\in\mathbb{C}:\operatorname{dist}(z,\gamma([0,1]))<\varepsilon\}$, so $\mathrm{Tube}_{\varepsilon}\subset\mathbb{C}\setminus\widetilde{\mathcal{P}}$. On the compact closure $\overline{\mathrm{Tube}}_{\varepsilon}$ the majorant $\tilde y(|\cdot|)$ is bounded: $M_{\varepsilon}:=\sup_{z\in\overline{\mathrm{Tube}}_{\varepsilon}}\tilde y(|z|)<\infty$. Define the Banach space
\begin{equation*}
\mathcal{X}_{\varepsilon}:=\Bigl\{h\ \text{holomorphic on }\mathrm{Tube}_{\varepsilon},\ \text{continuous on }\overline{\mathrm{Tube}}_{\varepsilon},\ h(0)=0,\ \sup_{z\in\overline{\mathrm{Tube}}_{\varepsilon}}|h(z)|\le M_{\varepsilon}+1\Bigr\}
\end{equation*}
with the supremum norm $\|\cdot\|_{\overline{\mathrm{Tube}}_{\varepsilon}}$, and the Volterra operator
\begin{equation}
(\Phi h)(z):=c_{0}\!\int_{[0,z]}\!K_{\mu}\,d\sigma
+c_{1}\!\int_{[0,z]}\!K_{\mu}\,h\,d\sigma
+c_{2}\!\int_{[0,z]}\!K_{\mu}\,h^{2}\,d\sigma,
\label{eq:volterra-tube}
\end{equation}
where $[0,z]$ denotes the portion of $\gamma$ from $0$ up to $z$ (extended continuously within $\mathrm{Tube}_{\varepsilon}$ by the straight line to $z$), and $K_{\mu}$ is the kernel of the Volterra equation \eqref{eq:volterra} in the variable $z$.

\begin{quote}
\textsc{Lemma B} (kernel integrability on the tube)\textsc{.} There exists $L_{\varepsilon}>0$ such that for every $z\in\overline{\mathrm{Tube}}_{\varepsilon}$,
\begin{equation*}
\int_{[0,z]}|K_{\mu}(z,\sigma)|\,|d\sigma|\le L_{\varepsilon}.
\end{equation*}
\end{quote}

\emph{Proof of Lemma B.} $|K_{\mu}(z,\sigma)|=|z^{1/\mu}-\sigma^{1/\mu}|^{\mu-1}|\sigma|^{1/\mu-1}/(\mu\Gamma(\mu))$. For $\mu\in(0,1]$, the exponent $\mu-1\in(-1,0]$ makes $(z^{1/\mu}-\sigma^{1/\mu})^{\mu-1}$ integrable in $\sigma$ near $\sigma=z$, and $\sigma^{1/\mu-1}$ is locally integrable near $\sigma=0$ since $1/\mu\ge 1$ for $\mu\le 1$, so $1/\mu-1\ge 0$. Uniform boundedness on $\overline{\mathrm{Tube}}_{\varepsilon}$ follows by compactness. \hfill$\square$

For $h_{1},h_{2}\in\mathcal{X}_{\varepsilon}$,
\begin{equation}
\|\Phi h_{1}-\Phi h_{2}\|_{\overline{\mathrm{Tube}}_{\varepsilon}}
\le L_{\varepsilon}\bigl(|c_{1}|+2|c_{2}|(M_{\varepsilon}+1)\bigr)\,\|h_{1}-h_{2}\|_{\overline{\mathrm{Tube}}_{\varepsilon}}.
\label{eq:contraction-bound}
\end{equation}
Rescale the tube by shrinking $\varepsilon$ if necessary (or subdivide $\gamma$ into finitely many subpaths, along each of which the analogous tube has sufficiently small $L_{\hat\varepsilon}$) so that $L_{\varepsilon}\bigl(|c_{1}|+2|c_{2}|(M_{\varepsilon}+1)\bigr)\le 1/2$. Then $\Phi$ is a contraction on $\mathcal{X}_{\varepsilon}$ and has a unique fixed point $g_{\varepsilon}$ agreeing with $g$ on $\{|z|<\tilde R\}\cap\mathrm{Tube}_{\varepsilon}$. By the identity principle, $g$ extends to a holomorphic function on $\mathrm{Tube}_{\varepsilon}$; and since $z_{0}\in\mathbb{C}\setminus\widetilde{\mathcal{P}}$ was arbitrary, $g$ extends holomorphically to all of $\mathbb{C}\setminus\widetilde{\mathcal{P}}$. The subdivision step is finite because $\gamma([0,1])$ is compact, $\tilde y$ is locally bounded away from its poles, and the kernel bound is uniform on compacts disjoint from $\widetilde{\mathcal{P}}$. Hence the extension is globally defined on $\mathbb{C}\setminus\widetilde{\mathcal{P}}$.

\emph{Step 5: exact residue at a finite pole.} By Step 4, $g$ is holomorphic on $\mathbb{C}\setminus\widetilde{\mathcal{P}}$; singularities of $g$, if any, lie in $\widetilde{\mathcal{P}}$ or in additional isolated points where the bootstrap fails. Let $z_{*}\in\mathbb{C}$ be an isolated singularity of $g$. We claim $z_{*}$ is a simple pole with residue $-(z_{*})^{(\mu-1)/\mu}/(c_{2}\mu)$. Near $z_{*}$, write $g(z)=A/(z-z_{*})+G(z)$ where $G$ is holomorphic at $z_{*}$, and substitute into the Volterra equation:
\begin{equation}
g(z)=c_{0}\!\int_{0}^{z}\!K_{\mu}(z,\sigma)\,d\sigma
+c_{1}\!\int_{0}^{z}\!K_{\mu}(z,\sigma)\Bigl[\frac{A}{\sigma-z_{*}}+G(\sigma)\Bigr]d\sigma
+c_{2}\!\int_{0}^{z}\!K_{\mu}(z,\sigma)\Bigl[\frac{A^{2}}{(\sigma-z_{*})^{2}}+\frac{2AG(\sigma)}{\sigma-z_{*}}+G(\sigma)^{2}\Bigr]d\sigma.
\label{eq:residue-substitution}
\end{equation}
As $z\to z_{*}$, the dominant singularities on both sides must match. The left-hand side has an order-$1$ pole. On the right-hand side, the term $c_{2}A^{2}\int_{0}^{z}K_{\mu}(z,\sigma)/(\sigma-z_{*})^{2}\,d\sigma$ develops --- by a Laplace-type analysis with the kernel $K_{\mu}(z,\sigma)\sim(z^{1/\mu}-\sigma^{1/\mu})^{\mu-1}\sigma^{1/\mu-1}/(\mu\Gamma(\mu))$ --- a pole of order $1$ at $z=z_{*}$ with leading behavior
\begin{equation}
c_{2}A^{2}\cdot\bigl[-\mu\,(z_{*})^{(1-\mu)/\mu}\bigr]\cdot\frac{1}{z-z_{*}},
\label{eq:residue-balance}
\end{equation}
obtained by the substitution $u=z^{1/\mu}-\sigma^{1/\mu}$ and evaluation of the resulting Mellin integral: the factor $-\mu(z_{*})^{(1-\mu)/\mu}$ is the Jacobian of the change of variables $\sigma=(z^{1/\mu}-u)^{\mu}$ at $u\to 0$, $\sigma\to z_{*}$, combined with the residue of $1/(\sigma-z_{*})^{2}$ under this change. Matching the $(z-z_{*})^{-1}$ coefficients,
\begin{equation}
A=c_{2}A^{2}\cdot\bigl[-\mu\,(z_{*})^{(1-\mu)/\mu}\bigr],
\label{eq:residue-match}
\end{equation}
which gives $A=0$ or
\begin{equation}
\Res(g,z_{*})=A=-\frac{1}{c_{2}\,\mu\,(z_{*})^{(1-\mu)/\mu}}=-\frac{(z_{*})^{(\mu-1)/\mu}}{c_{2}\,\mu},
\label{eq:residue-value}
\end{equation}
which is \eqref{eq:residue}. This supersedes both the ``$-\Gamma(1-\mu)/c_{2}$'' of one earlier draft and the ``$-1/c_{2}$'' of another: the correct residue is $z_{*}$-\emph{dependent} through the Jacobian factor. At $\mu=1$ it reduces to $\Res=-1/c_{2}$, matching the classical Riccati; the $\Gamma(1-\mu)$ factors of earlier versions are absorbed into the kernel Jacobian during the Mellin evaluation. Higher-order poles are excluded: if $g$ had a pole of order $\ge 2$ at $z_{*}$, the left-hand side would contribute a $(z-z_{*})^{-2}$ singularity, while the right-hand side (through the same Jacobian analysis) contributes at most $(z-z_{*})^{-1}$, a contradiction.

\emph{Step 6: non-accumulation via a uniform residue lower bound.} Suppose, for contradiction, that $\mathcal{P}_{g}$ has a finite accumulation point $z_{\infty}\in\mathbb{C}$, and let $\{z_{*,n}\}\subset\mathcal{P}_{g}$ be a sequence with $z_{*,n}\to z_{\infty}$. Choose $\delta>0$ with $\overline{D}(z_{\infty},2\delta)\cap\widetilde{\mathcal{P}}=\varnothing$, possible since $\widetilde{\mathcal{P}}$ is discrete. By Step 5, each $z_{*,n}$ has residue $A_{n}=-(z_{*,n})^{(\mu-1)/\mu}/(c_{2}\mu)$. For $n$ large, $z_{*,n}\in D(z_{\infty},\delta)$, so $|z_{*,n}|\in[|z_{\infty}|-\delta,\,|z_{\infty}|+\delta]$ is bounded away from $0$ (the case $z_{\infty}=0$ is excluded because $g$ is holomorphic at $0$ with $g(0)=0$). Hence
\begin{equation}
|A_{n}|\ge\frac{(|z_{\infty}|-\delta)^{(\mu-1)/\mu}}{|c_{2}|\,\mu}=:A_{\min}>0.
\label{eq:residue-lower}
\end{equation}
Let $\varepsilon_{n}:=\tfrac{1}{2}\operatorname{dist}\bigl(z_{*,n},\mathcal{P}_{g}\setminus\{z_{*,n}\}\bigr)$. Since $\{z_{*,n}\}$ accumulates at $z_{\infty}$, eventually $\varepsilon_{n}\to 0$. On the punctured disk $D(z_{*,n},\varepsilon_{n})\setminus\{z_{*,n}\}$,
\begin{equation}
|g(z)|\ge\frac{|A_{n}|}{|z-z_{*,n}|}-\sup_{z\in D(z_{*,n},\varepsilon_{n})}|G_{n}(z)|,
\label{eq:blowup-lower}
\end{equation}
where $G_{n}$ is holomorphic on $D(z_{*,n},2\varepsilon_{n})$. In particular, on the circle $|z-z_{*,n}|=\varepsilon_{n}/2$,
\begin{equation}
|g(z)|\ge\frac{2|A_{n}|}{\varepsilon_{n}}-\sup|G_{n}|\;\longrightarrow\;\infty
\qquad\text{as }n\to\infty.
\label{eq:blowup-circle}
\end{equation}
But by Step 3 extended over a tube covering $D(z_{\infty},\delta)$, we have $|g(z)|\le M:=\sup_{|\zeta|\le|z_{\infty}|+\delta}\tilde y(|\zeta|)<\infty$ on $\overline{D}(z_{\infty},\delta)\setminus\mathcal{P}_{g}$, since $\overline{D}(z_{\infty},\delta)$ is disjoint from $\widetilde{\mathcal{P}}$. This contradicts \eqref{eq:blowup-circle}. Hence $\mathcal{P}_{g}$ has no finite accumulation point.

Combining Steps 1--6, $g$ is meromorphic on $\mathbb{C}$ with discrete simple pole set $\mathcal{P}_{g}$ and residues given exactly by \eqref{eq:residue}.
\end{proof}

\begin{remark}[Summary of corrections versus earlier versions]
\label{rem:meromorphy-corrections}
\begin{enumerate}[label=(\alph*)]
\item The constant in the Gautschi-type bound is \emph{not} $1$ but $C_{\mu}\le 21/20$; Lemma A gives the sharp value $C_{\mu}=\Gamma(\mu_{*}+1)/\Gamma(2\mu_{*}+1)\approx 1.0428$ at $\mu_{*}\approx 0.1339$. All earlier claims using ``$\gamma_{k}\le 1$'' are hereby corrected.
\item The majorant ODE \eqref{eq:majorant-ode} carries the $C_{\mu}$ factor in front of \emph{every} nonlinear term, not just some of them. The explicit solution \eqref{eq:majorant-exp} is written with the correctly scaled $\lambda_{\pm}$.
\item The residue is $z_{*}$-\emph{dependent}: $\Res(g,z_{*})=-(z_{*})^{(\mu-1)/\mu}/(c_{2}\mu)$. At $\mu=1$ this reduces to $-1/c_{2}$, matching the classical Cole--Hopf. At $\mu<1$ the residue carries a nontrivial power of $z_{*}$ inherited from the kernel Jacobian.
\item Analytic continuation is by \emph{finitely many} tube contractions, each with a quantitatively bounded operator norm (Lemma B together with the estimate \eqref{eq:contraction-bound}). The subdivision uses compactness of $\gamma$.
\item Non-accumulation is proved by the \emph{uniform lower bound} \eqref{eq:residue-lower}--\eqref{eq:blowup-circle} on the residues, not by an unquantified ``bounded on a deleted neighborhood'' appeal. The uniform bound on $|A_{n}|$ is essential and was missing earlier.
\item The simple-pole structure is proved in \emph{both} directions: existence of a simple pole by the matched-asymptotic balance \eqref{eq:residue-balance}, and exclusion of higher-order poles by the impossibility of matching $(z-z_{*})^{-k}$ for $k\ge 2$ on the right-hand side.
\end{enumerate}
No use is made of a Leibniz rule or of a Cole--Hopf linearization anywhere in the proof.
\end{remark}

\subsection{Global convergence on the positive real axis}
\label{sec:heston-global}

\begin{lemma}[Absence of poles on the positive real axis]
\label{lem:no-real-poles}
Let $y(t)$ be the solution of \eqref{eq:fractional-riccati} on $[0,\infty)$. Then $g(z)=y(z^{1/\mu})$ has no poles on $[0,\infty)$. Consequently, for every $T>0$,
\begin{equation}
\delta_{T}:=\mathrm{dist}\bigl([0,T^{\mu}],\mathcal{P}_{g}\bigr)>0.
\label{eq:delta-T}
\end{equation}
\end{lemma}

\begin{proof}
The Volterra equation \eqref{eq:volterra} has a unique continuous solution on $[0,\infty)$ for the parameter ranges arising in rough Heston (the quadratic nonlinearity is locally Lipschitz and the kernel is integrable). A pole of $g$ at $z=t^{\mu}>0$ would correspond to a blow-up of $y(t)$ at finite positive $t$, contradicting the global existence on $[0,\infty)$. Hence $\mathcal{P}_{g}\cap[0,\infty)=\varnothing$, and since $\mathcal{P}_{g}$ is discrete, the distance $\delta_{T}$ is strictly positive.
\end{proof}

\begin{lemma}[Gon\v{c}ar condition on compact intervals]
\label{lem:goncar-interval}
For every $T>0$, the diagonal Pad\'e approximants $[M/M]_{g}=R_{M}$ satisfy the Gon\v{c}ar condition on $K=[0,T^{\mu}]$: the number of poles of $R_{M}$ in $K$ is uniformly bounded in $M$.
\end{lemma}

\begin{proof}
By Theorem~\ref{prop:meromorphy}, $g$ is meromorphic with discrete pole set $\mathcal{P}_{g}$. Let $\delta_{T}>0$ be as in \eqref{eq:delta-T}. The poles of $R_{M}$ fall into two classes:
\begin{enumerate}[label=(\alph*)]
\item \emph{True-pole approximants.} These converge to the elements of $\mathcal{P}_{g}$ (de~Montessus de~Ballore). For $M$ sufficiently large, every such pole lies within $\delta_{T}/2$ of some element of $\mathcal{P}_{g}$, hence outside $K$.
\item \emph{Froissart doublets.} By Proposition~\ref{prop:doublet-count}, the number of Froissart doublets of $R_{M}$ in any compact subset of $\mathbb{C}\setminus\mathcal{P}_{g}$ is $O(1)$ as $M\to\infty$.
\end{enumerate}
Therefore, for all sufficiently large $M$, the only poles of $R_{M}$ that can lie in $K$ are Froissart doublets, whose number is uniformly bounded. The Gon\v{c}ar condition holds.
\end{proof}

\begin{theorem}[Global Pad\'e--M\"untz representation]
\label{thm:global}
Let $y(t)$ solve the fractional Riccati equation \eqref{eq:fractional-riccati} with $\mu\in(0,1]$. Then for every $t\ge 0$,
\begin{equation}
\boxed{\;y(t)=\lim_{M\to\infty}R_{M}(t^{\mu})=\lim_{M\to\infty}\frac{P_{M}(t^{\mu})}{Q_{M}(t^{\mu})}.\;}
\label{eq:global-pade}
\end{equation}
Moreover, for every $T>0$, the convergence is uniform on $[0,T]$.
\end{theorem}

\begin{proof}
By Theorem~\ref{prop:meromorphy}, $g(z)=y(z^{1/\mu})$ is meromorphic on $\mathbb{C}$. By Lemma~\ref{lem:no-real-poles}, $[0,T^{\mu}]$ is disjoint from the pole set $\mathcal{P}_{g}$. Lemma~\ref{lem:goncar-interval} establishes the Gon\v{c}ar condition on $[0,T^{\mu}]$. Applying Gon\v{c}ar's uniform convergence theorem (Theorem~\ref{thm:goncar}) to the compact $K=[0,T^{\mu}]$, we obtain $R_{M}\to g$ uniformly on $[0,T^{\mu}]$. Since $T>0$ is arbitrary, uniform convergence holds on every compact subset of $[0,\infty)$, and pointwise convergence holds for every $t\ge 0$.
\end{proof}

\subsection{The bordered rank-one update and incremental Hankel solver}
\label{sec:rank-one}

The passage from Pad\'e order $M$ to order $M+1$ requires the solution of the augmented Hankel system $H_{M+1}\bm{q}_{M+1}=-\bm{b}_{M+1}$. Direct inversion costs $O(M^{3})$. The Hankel structure permits an incremental $O(M^{2})$ update via a bordered rank-one correction.

\begin{lemma}[Bordered structure of the augmented Hankel system]
\label{lem:bordered}
Let $P_{M}$ be the cyclic column permutation that moves the first column to the last position. Then
\begin{equation}
\hat{H}_{M+1}:=H_{M+1}P_{M}=
\begin{pmatrix}
H_{M} & \mathbf{u}_{M}\\[2pt]
\mathbf{v}_{M}^{T} & d_{M}
\end{pmatrix},
\label{eq:bordered}
\end{equation}
where
\begin{align}
\mathbf{u}_{M}&=(a_{M+1},a_{M+2},\ldots,a_{2M})^{T}\in\mathbb{C}^{M},\\
\mathbf{v}_{M}^{T}&=(a_{2M},a_{2M-1},\ldots,a_{M+1})\in\mathbb{C}^{1\times M},\\
d_{M}&=a_{2M+1}\in\mathbb{C}.
\end{align}
\end{lemma}

\begin{proof}
The matrix $H_{M+1}P_{M}$ has columns $2,3,\ldots,M+1,1$ of $H_{M+1}$. The first $M$ columns are therefore the top-right $M\times M$ block of $H_{M+1}$, which is $H_{M}$ by direct inspection of \eqref{eq:hankel-system}. The last column is the first column of $H_{M+1}$, namely $(a_{M+1},\ldots,a_{2M+1})^{T}$, giving $\mathbf{u}_{M}$ and $d_{M}$. The last row of the first $M$ columns is $(a_{2M},a_{2M-1},\ldots,a_{M+1})$, giving $\mathbf{v}_{M}^{T}$.
\end{proof}

\begin{theorem}[Bordered rank-one update of the inverse]
\label{thm:rank-one}
Assume $H_{M}$ is invertible and define the Schur complement
\begin{equation}
s_{M}=d_{M}-\mathbf{v}_{M}^{T}H_{M}^{-1}\mathbf{u}_{M}.
\label{eq:schur}
\end{equation}
If $s_{M}\neq 0$, then $H_{M+1}$ is invertible and
\begin{equation}
\hat{H}_{M+1}^{-1}=
\begin{pmatrix}
H_{M}^{-1}+\dfrac{1}{s_{M}}\mathbf{z}_{M}\mathbf{w}_{M}^{T} & -\dfrac{1}{s_{M}}\mathbf{z}_{M}\\[8pt]
-\dfrac{1}{s_{M}}\mathbf{w}_{M}^{T} & \dfrac{1}{s_{M}}
\end{pmatrix},
\label{eq:rank-one-update}
\end{equation}
where
\begin{equation}
\mathbf{z}_{M}=H_{M}^{-1}\mathbf{u}_{M},\qquad
\mathbf{w}_{M}^{T}=\mathbf{v}_{M}^{T}H_{M}^{-1}.
\label{eq:zw}
\end{equation}
The update from $H_{M}^{-1}$ to $\hat{H}_{M+1}^{-1}$ requires two matrix-vector products, one inner product, and one rank-one outer product, for a total cost of $O(M^{2})$ arithmetic operations.
\end{theorem}

\begin{proof}
This is the standard block inversion formula for a $2\times 2$ block matrix with invertible leading block and nonzero Schur complement. The matrix $\mathbf{z}_{M}\mathbf{w}_{M}^{T}$ is rank one, hence the name. The cost breakdown is: $H_{M}^{-1}\mathbf{u}_{M}$ costs $O(M^{2})$; $\mathbf{v}_{M}^{T}H_{M}^{-1}$ costs $O(M^{2})$; $\mathbf{v}_{M}^{T}\mathbf{z}_{M}$ costs $O(M)$; the outer product $\mathbf{z}_{M}\mathbf{w}_{M}^{T}$ costs $O(M^{2})$.
\end{proof}

\begin{corollary}[Incremental denominator solver]
\label{cor:incremental-solver}
Given $H_{M}^{-1}$ and the solution $\mathbf{q}_{M}$ of $H_{M}\mathbf{q}_{M}=-\mathbf{b}_{M}$, the augmented solution $\mathbf{q}_{M+1}$ of $H_{M+1}\mathbf{q}_{M+1}=-\mathbf{b}_{M+1}$ is obtained as follows:
\begin{enumerate}[label=(\roman*)]
\item Compute $\mathbf{z}_{M}=H_{M}^{-1}\mathbf{u}_{M}$ and $\mathbf{w}_{M}^{T}=\mathbf{v}_{M}^{T}H_{M}^{-1}$ via \eqref{eq:zw}.
\item Compute the Schur complement $s_{M}=d_{M}-\mathbf{w}_{M}^{T}\mathbf{u}_{M}$.
\item Update the inverse via \eqref{eq:rank-one-update} to obtain $\hat{H}_{M+1}^{-1}$.
\item Solve $\hat{H}_{M+1}\hat{\mathbf{q}}_{M+1}=-\mathbf{b}_{M+1}$ using the updated inverse.
\item Permute back: $\mathbf{q}_{M+1}=P_{M}\hat{\mathbf{q}}_{M+1}$.
\end{enumerate}
The total cost of the increment is $O(M^{2})$. Amortized over $M=1,2,\ldots,N$, the cumulative cost to reach order $N$ is $O(N^{2})$.
\end{corollary}

\begin{remark}[Connection to the Jacobi truncation]
\label{rem:jacobi-update}
Theorem~\ref{thm:rank-one} is the algebraic realization of the Jacobi truncation described in Theorem~\ref{thm:hankel-jacobi}. Adding one row and column to $J_{M}$ to produce $J_{M+1}$ corresponds to bordering the Hankel matrix. The rank-one correction $\mathbf{z}_{M}\mathbf{w}_{M}^{T}$ encodes the spectral perturbation induced by the new Jacobi entry, and the resulting Pad\'e iterate $R_{M+1}$ is the resolvent of the enlarged operator at the basis vector $e_{0}$.
\end{remark}

\subsection{Spectral interpretation and operator-theoretic synthesis}
\label{sec:heston-spectral}

\begin{theorem}[Hankel solver as Jacobi truncation]
\label{thm:hankel-jacobi}
The incremental Hankel solver of Corollary~\ref{cor:incremental-solver}, which builds $H_{M+1}^{-1}$ from $H_{M}^{-1}$ via the bordered rank-one update \eqref{eq:rank-one-update} in $O(M^{2})$ operations, is the algebraic realization of the Jacobi truncation: it adds one row and column to $J_{M}$ to produce $J_{M+1}$, and the resulting Pad\'e iterate $R_{M+1}$ is the resolvent of the enlarged operator at the basis vector $e_{0}$.
\end{theorem}

\begin{corollary}[Convergence certificate as spectral gap]
\label{cor:spectral-gap}
The successive-difference bound \eqref{eq:error-final} certifies that the iterate $R_{M}$ has entered the asymptotic regime where the tail is geometrically controlled by $\exp(-Mg(\xi,\infty))$, with $g$ the Green's function of the cut plane in the variable $\xi=z=t^{\mu}$. The threshold $|\Delta_{M}|^{2}/(|\Delta_{M-1}|-|\Delta_{M}|)<2^{-\mathrm{bits}}$ detects the spectral gap between the working point and the Stahl compact.
\end{corollary}

\begin{theorem}[Compiled spectral reconstruction]
\label{thm:compiled}
Let $\Phi(u,T)$ denote the cumulant generating function of the rough Heston model under standard normalizations:
\begin{equation}
\Phi(u,T)=\theta\lambda\int_{0}^{T}h(u,t)\,dt+V_{0}I^{1-\mu}[h(u,\cdot)](T),
\label{eq:cgf}
\end{equation}
where $h$ solves the fractional Riccati equation \eqref{eq:fractional-riccati} with parameters $(c_{0},c_{1},c_{2})$ depending on $u$. Write $\xi=T^{\mu}$. Then $\Phi$ admits the representation
\begin{equation}
\Phi(u,T)=T\sum_{k=0}^{\infty}d_{k}(u)\xi^{k},
\label{eq:cgf-series}
\end{equation}
where
\begin{align}
d_{0}(u)&=V_{0}c_{0}(u),\\
d_{k}(u)&=\frac{\theta\lambda a_{k}(u)}{k\mu+1}+\frac{V_{0}\Gamma((k+1)\mu+1)}{\Gamma(k\mu+2)}a_{k+1}(u),\qquad k\ge 1,
\label{eq:dk}
\end{align}
and the coefficients $a_{k}(u)$ are the Puiseux coefficients of Theorem~\ref{thm:riccati-series}. The series \eqref{eq:cgf-series} is a power series in $\xi=T^{\mu}$ and is resummed via diagonal Pad\'e in $\xi$:
\begin{equation}
\Phi(u,T)=T\cdot\lim_{M\to\infty}[M/M]_{S(u,\cdot)}(\xi),\qquad S(u,\xi)=\sum_{k=0}^{\infty}d_{k}(u)\xi^{k}.
\label{eq:compiled}
\end{equation}
Convergence is governed by Theorem~\ref{thm:global} applied to the meromorphic function $S(u,\cdot)$, and is certified a-posteriori by Corollary~\ref{cor:practical-cert}.
\end{theorem}

\begin{proof}
Substituting the M\"untz series $h(u,t)=\sum_{k=1}^{\infty}a_{k}(u)t^{k\mu}$ into \eqref{eq:cgf} and using \eqref{eq:Ir-power} for the fractional integral term yields
\[
\Phi(u,T)=\theta\lambda\sum_{k=1}^{\infty}a_{k}(u)\frac{T^{k\mu+1}}{k\mu+1}+V_{0}\sum_{k=1}^{\infty}a_{k}(u)\frac{\Gamma(k\mu+1)}{\Gamma((k-1)\mu+2)}T^{(k-1)\mu+1}.
\]
Factoring out $T$ and reindexing the second sum with $j=k-1$ gives
\[
\Phi(u,T)=T\Biggl[V_{0}a_{1}(u)\Gamma(\mu+1)+\sum_{k=1}^{\infty}\biggl(\frac{\theta\lambda a_{k}(u)}{k\mu+1}+\frac{V_{0}\Gamma((k+1)\mu+1)}{\Gamma(k\mu+2)}a_{k+1}(u)\biggr)\xi^{k}\Biggr].
\]
Since $a_{1}(u)=c_{0}(u)/\Gamma(\mu+1)$ and $\Gamma(2)=1$, the constant term simplifies to $V_{0}c_{0}(u)=d_{0}(u)$. The coefficients of $\xi^{k}$ for $k\ge 1$ are exactly $d_{k}(u)$, proving \eqref{eq:cgf-series}. The Pad\'e resummation \eqref{eq:compiled} follows by applying the construction of Section~\ref{sec:heston-pade} to the series $S(u,\xi)$.
\end{proof}

\subsection{Algorithmic summary}
\label{sec:heston-algorithm}

Given $(c_{0},c_{1},c_{2},\mu)$ and Pad\'e order $M$:
\begin{enumerate}
\item Compute the Puiseux coefficients $a_{1},\dots,a_{2M}$ by the recurrence \eqref{eq:a1}--\eqref{eq:ak}. Cost: $O(M^{2})$.
\item Form the Hankel matrix $H_{M}$ and right-hand side $\bm{b}$ from \eqref{eq:hankel-system}.
\item Solve $H_{M}\bm{q}=-\bm{b}$ for the denominator coefficients. Use direct factorization for the base case $M=1$, then increment via the bordered rank-one update of Corollary~\ref{cor:incremental-solver} for each subsequent order. Cumulative cost to order $M$: $O(M^{2})$.
\item Recover the numerator coefficients via \eqref{eq:p-eqs}.
\item Evaluate $y_{M}(t)=P_{M}(t^{\mu})/Q_{M}(t^{\mu})$ at any $t\ge 0$. Cost: $O(M)$ per query.
\end{enumerate}
Steps 1--4 are independent of $t$. The entire family $\{y_{M}(t):t\ge 0\}$ is encoded in a single rational function in $z=t^{\mu}$. There is no grid, no time stepping, no Newton iteration, and no fixed-point loop at evaluation time.

\subsection{Geometry of the Stahl compact for meromorphic Mittag--Leffler functions}
\label{sec:heston-stahl-geometry}

The function $g(z)=y(z^{1/\mu})$ of Proposition~\ref{prop:meromorphy} is meromorphic on $\mathbb{C}$ with discrete pole set $\mathcal{P}_{g}$, a ratio of entries of a fundamental matrix of Mittag--Leffler functions. This subsection describes the Stahl compact $\Delta_{g}$ of such a function: which compacts are admissible, what shape the extremal arcs take, and why none of that geometry enters the algorithm of Section~\ref{sec:heston-algorithm}.

\begin{theorem}[Arcs connect poles to infinity]
\label{thm:arcs-to-infinity}
Let $g$ be meromorphic on $\mathbb{C}$ with discrete pole set $\mathcal{P}_{g}$. The admissible compacts of Definition~\ref{def:admissible} are precisely the compacts containing $\mathcal{P}_{g}$, and the Stahl compact $\Delta_{g}$ minimizes logarithmic capacity among them. Since $\mathcal{P}_{g}$ is discrete, the minimal capacity is achieved by adjoining to the poles a system of arcs $\gamma$ connecting the poles to each other and to infinity, so that
\begin{equation}
\Delta_{g}=\overline{\mathcal{P}_{g}}\cup\gamma
\label{eq:stahl-compact-structure}
\end{equation}
is a connected set of minimal logarithmic capacity.
\end{theorem}

\begin{proof}
A compact $K$ is admissible for $g$ if and only if $g$ admits single-valued meromorphic continuation on $\overline{\mathbb{C}}\setminus K$ (Definition~\ref{def:admissible}). Since $g$ is already single-valued and meromorphic on all of $\mathbb{C}$, the continuation exists on $\overline{\mathbb{C}}\setminus K$ exactly when $K$ contains every pole of $g$; removing a pole from $K$ leaves a point of $\overline{\mathbb{C}}\setminus K$ at which $g$ has no meromorphic extension only if that point is an essential singularity, which is excluded, so the condition reduces to $\mathcal{P}_{g}\subset K$. Among admissible compacts, Theorem~\ref{thm:stahl-geometric} selects the unique $\Delta_{g}$ of minimal capacity. The connectivity of the extremal configuration and the presence of connecting arcs follow from the variational characterization of Remark~\ref{rem:variational}: within the homotopy class of admissible compacts, the capacity minimum is attained on a connected set, and connecting a discrete point set into a single continuum of minimal capacity requires arcs joining the poles to each other and, when $\mathcal{P}_{g}$ is infinite, to infinity.
\end{proof}

\begin{theorem}[The arcs are not straight]
\label{thm:arcs-not-straight}
For generic parameters $(c_{0},c_{1},c_{2})$ with $c_{2}\ne 0$, the arcs of $\gamma$ are not straight line segments. They are analytic curves satisfying the $S$-property
\begin{equation}
\frac{\partial g_{\Delta_{g}}}{\partial n_{+}}(s)=\frac{\partial g_{\Delta_{g}}}{\partial n_{-}}(s),
\qquad s\in\gamma\setminus\{\text{endpoints}\},
\label{eq:s-property-arcs}
\end{equation}
where $g_{\Delta_{g}}(\cdot,\infty)$ is the Green's function of $\overline{\mathbb{C}}\setminus\Delta_{g}$ with pole at infinity and $n_{\pm}$ are the two unit normals. The curvature of each arc is determined by the local distribution of the poles and residues of $g$. A straight line segment satisfies the $S$-property only when the harmonic measure is symmetric about the line, which requires a symmetry of the pole--residue configuration of $g$ about that line and is nongeneric in $(c_{0},c_{1},c_{2})$.
\end{theorem}

\begin{proof}
The arcs of any Stahl compact are trajectories of a quadratic differential and hence analytic curves satisfying \eqref{eq:s-property-arcs} (Theorem~\ref{thm:s-property}). The $S$-property equates the two normal derivatives of the Green's function along each arc, and the Green's function is the logarithmic potential of the equilibrium measure of $\Delta_{g}$ up to the Robin constant; its normal derivatives at a point $s\in\gamma$ are therefore functionals of the full configuration $\Delta_{g}$, which for the Mittag--Leffler ratio $g$ is determined by the poles and residues of $g$. If an arc were a straight segment, reflection across the supporting line would have to preserve the two one-sided harmonic measures, forcing the pole--residue configuration of $g$ to be symmetric about that line. The zeros of the Mittag--Leffler entries of the fundamental matrix depend analytically and nontrivially on $(c_{0},c_{1},c_{2})$, and the symmetry condition imposes nontrivial analytic constraints on these parameters; the constrained set is a proper analytic subvariety of the parameter space, so for generic $(c_{0},c_{1},c_{2})$ with $c_{2}\ne 0$ no arc is straight.
\end{proof}

\begin{theorem}[Algorithmic irrelevance of the arc geometry]
\label{thm:arc-irrelevance}
The incremental Hankel solver of Corollary~\ref{cor:incremental-solver}, the computable error bound of Theorem~\ref{thm:error-bound}, and the uniform convergence theorem (Theorem~\ref{thm:global}) do not require knowledge of the arc geometry of $\gamma$. The Pad\'e approximant $R_{M}$ is constructed from the M\"untz coefficients $a_{k}$ alone; the Riemann--Hilbert matrix $Y$ is built from the orthogonal polynomials $Q_{n}$ for the discrete measure determined by the poles and residues of $g$; and the convergence certificate depends only on the successive Pad\'e differences $\Delta_{k}$ of Definition~\ref{def:differences}. The arcs $\gamma$ are the theoretical backbone of the convergence theory, not a computational input.
\end{theorem}

\begin{proof}
Each claim is read off from the corresponding construction. The solver of Corollary~\ref{cor:incremental-solver} consumes only the Hankel data $a_{1},\dots,a_{2M}$ produced by the recurrence \eqref{eq:a1}--\eqref{eq:ak}; no geometric quantity enters. The error bound \eqref{eq:error-uniform} of Theorem~\ref{thm:error-bound} is stated entirely in terms of the observable differences $\Delta_{k}(z)$ and their contraction ratio. The proof of Theorem~\ref{thm:global} invokes the Stahl compact only through the qualitative facts that $\mathcal{P}_{g}$ is discrete and avoids $[0,\infty)$; the shape of the arcs never appears. The arc geometry governs the exponential convergence \emph{rate} through the Green's function $g_{\Delta_{g}}(z,\infty)$ (Corollary~\ref{cor:exponential-rate}), but the algorithm certifies its accuracy a~posteriori without evaluating that rate.
\end{proof}

\begin{remark}[The arcs organize the Pad\'e poles]
\label{rem:arcs-organize-poles}
The arcs are the structural framework onto which the Pad\'e approximant organizes its poles. At order $n$, the $n$ poles of $[n/n]_{g}$ are distributed near the true poles of $g$ and along the arcs of $\gamma$. As $n\to\infty$ the poles condense: the true poles are approximated by nearby Pad\'e poles (de~Montessus de~Ballore), while the remaining Pad\'e poles distribute along $\gamma$ according to the equilibrium measure $\mu_{\Delta_{g}}$ (Theorem~\ref{thm:zeros-equilibrium}). Froissart doublets (Definition~\ref{def:froissart}) are the mechanism by which spurious poles are canceled by nearby zeros, mimicking the residue structure of the true spectral measure.
\end{remark}

\begin{remark}[Accumulation at infinity]
\label{rem:accumulation-infinity}
The pole set $\mathcal{P}_{g}$ may accumulate at infinity, as the zeros of the Mittag--Leffler entries do. In that case the arcs connect the finite poles to each other and to infinity in a tree structure. The Stahl compact is then unbounded in $\mathbb{C}$ (compact in $\overline{\mathbb{C}}$), but its intersection with any bounded set contains only finitely many poles and arc segments. The equilibrium measure $\mu_{\Delta_{g}}$ is still well-defined, and the convergence theory of Section~\ref{sec:heston-global} applies without modification.
\end{remark}

\section{The essential difference}
\label{sec:conclusion}

Taylor approximation asks: \emph{what polynomial has the same derivatives at this point?} It is local, and its radius is the distance to the nearest singularity.

Pad\'e approximation asks: \emph{what spectral measure produces these moments, and what is its Stieltjes transform?} It is global, and its convergence domain is the maximal region on which the Riemann--Hilbert factorization exists. The poles are not defects; they are the discrete spectrum of a truncated Jacobi operator, and their condensation on the equilibrium contour is the signature of the underlying Riemann surface. The approximant continues $f$ because it is solving for the measure that generated the series, and that reconstruction is valid wherever the logarithmic potential of the equilibrium measure controls the error.

In the rough Heston application, this distinction is decisive. The M\"untz--Tau substitution produces a Puiseux series whose radius of convergence is finite and generically much smaller than the time horizon of practical interest. The series alone is therefore a local representation, useful only on $[0,T_{0}]$ for the Banach radius $T_{0}$ of Theorem~\ref{thm:radius}. Pad\'e resummation in the variable $z=t^{\mu}$ converts the local representation into a global one, valid on all of $[0,\infty)$ by the meromorphy established in Theorem~\ref{prop:meromorphy}, with computable a-posteriori error control. The mechanism is the same in every section of this document: the discrete poles of the truncation reorganize themselves into the equilibrium configuration of the function's true singular set, and the rational function so produced is the Stieltjes transform of the spectral measure that the original series was secretly encoding.

\begin{thebibliography}{99}

\bibitem{Stahl1986a}
H.~Stahl, \emph{Orthogonal polynomials with complex-valued weight functions. I}, Constr.\ Approx.\ \textbf{2} (1986), 225--240.

\bibitem{Stahl1986b}
H.~Stahl, \emph{Orthogonal polynomials with complex-valued weight functions. II}, Constr.\ Approx.\ \textbf{2} (1986), 241--251.

\bibitem{Stahl1997}
H.~Stahl, \emph{The convergence of Pad\'e approximants to functions with branch points}, J.~Approx.\ Theory \textbf{91} (1997), 139--204.

\bibitem{FokasItsKitaev1992}
A.~S.~Fokas, A.~R.~Its, and A.~V.~Kitaev, \emph{The isomonodromy approach to matrix models in 2D quantum gravity}, Comm.\ Math.\ Phys.\ \textbf{147} (1992), 395--430.

\bibitem{DeiftZhou1993}
P.~Deift and X.~Zhou, \emph{A steepest descent method for oscillatory Riemann--Hilbert problems}, Ann.\ of Math.\ \textbf{137} (1993), 295--368.

\bibitem{Deift1999}
P.~Deift, \emph{Orthogonal Polynomials and Random Matrices: A Riemann--Hilbert Approach}, Courant Lecture Notes 3, AMS, 1999.

\bibitem{Goncar1983}
A.~A.~Gon\v{c}ar, \emph{On the convergence of Pad\'e approximants for some classes of meromorphic functions}, Math.\ USSR Sb.\ \textbf{26} (1983), 551--575.

\bibitem{BakerGravesMorris1996}
G.~A.~Baker and P.~Graves-Morris, \emph{Pad\'e Approximants}, 2nd ed., Cambridge University Press, 1996.

\bibitem{Montessus1902}
R.~de~Montessus de~Ballore, \emph{Sur les fractions continues alg\'ebriques}, Bull.\ Soc.\ Math.\ France \textbf{30} (1902), 28--36.

\bibitem{Nuttall1970}
J.~Nuttall, \emph{The convergence of Pad\'e approximants of meromorphic functions}, J.~Math.\ Anal.\ Appl.\ \textbf{31} (1970), 147--153.

\bibitem{Pommerenke1973}
Ch.~Pommerenke, \emph{Pad\'e approximants and convergence in capacity}, J.~Math.\ Anal.\ Appl.\ \textbf{41} (1973), 775--780.

\bibitem{Suetin2021}
S.~P.~Suetin, \emph{On an example of the approximation of a function by the diagonal Pad\'e approximants}, Proc.\ Steklov Inst.\ Math.\ \textbf{315} (2021), 241--255.

\bibitem{Esmaeili2015}
S.~Esmaeili and M.~Shamsi, \emph{The M\"untz--Legendre Tau method for fractional differential equations}, Appl.\ Math.\ Modelling \textbf{40} (2015), 671--684.

\bibitem{Ecalle1981}
J.~\'Ecalle, \emph{Les fonctions r\'esurgentes}, Vols.~I--III, Publ.\ Math.\ d'Orsay, 1981--1985.

\bibitem{Costin2008}
O.~Costin, \emph{Asymptotics and Borel Summability}, Chapman \& Hall/CRC, 2008.

\bibitem{ElEuchRosenbaum2019}
O.~El~Euch and M.~Rosenbaum, \emph{The characteristic function of rough Heston models}, Math.\ Finance \textbf{29} (2019), 3--38.

\bibitem{GatheralRadoicic2019}
J.~Gatheral and B.~Radoi\v{c}i\'c, \emph{Rational approximation of the rough Heston solution}, Int.\ J.~Theor.\ Appl.\ Finance \textbf{22} (2019), 1950010.

\bibitem{CallegaroGrasselliPages2021}
G.~Callegaro, M.~Grasselli, and G.~Pag\`es, \emph{Fast hybrid schemes for fractional Riccati equations (rough is not so tough)}, Math.\ Oper.\ Res.\ \textbf{46} (2021), 221--254.

\end{thebibliography}

\end{document}
